% FC 2027 paper — The Cost of Custody
% Target: 15pp Springer LNCS main body, unlimited references + appendix
% Double-blind submission

\documentclass[runningheads]{llncs}

% ===================== Packages =====================
\usepackage[T1]{fontenc}
\usepackage[utf8]{inputenc}
\usepackage{graphicx}
\usepackage{amsmath,amssymb}
\usepackage{mathtools}
\usepackage{booktabs}
\usepackage{array}
\usepackage{tabularx}
\usepackage{multirow}
\usepackage[inline]{enumitem}
\usepackage{xcolor}
\usepackage{tikz}
\usetikzlibrary{shapes,arrows.meta,positioning,calc,fit,backgrounds}
\usepackage{pgfplots}
\pgfplotsset{compat=1.18}
\usepackage{siunitx}
\usepackage{url}
\usepackage{microtype}
\usepackage[bookmarksnumbered,hidelinks]{hyperref}

% ===================== Custom macros =====================
% Transaction size units
\newcommand{\vB}{\text{\,vB}}
\newcommand{\satvB}{\text{\,sat/vB}}

% Design / proposal names — roman, short OP_ prefix throughout. The
% full opcode names are glossed once on first use; thereafter we use
% these shorthands uniformly. For an actual opcode invocation in a
% code-like context (e.g.\ in an assertion or script listing), wrap
% locally in \texttt{...}.
\newcommand{\opctv}{OP\_CTV}
\newcommand{\opccv}{OP\_CCV}
\newcommand{\opvault}{OP\_VAULT}
\newcommand{\opvaultrecover}{OP\_VAULT\_RECOVER}
\newcommand{\opcat}{OP\_CAT}
\newcommand{\opcsfs}{OP\_CSFS}
\newcommand{\catcsfs}{\opcat{}+\opcsfs{}}

% Bitcoin concept names
\newcommand{\csv}{\texttt{OP\_CSV}}
\newcommand{\sighash}[1]{\texttt{SIGHASH\_#1}}

% Design-space axis labels --- canonical, see conference/AXES.md.
% Use these macros in every axis reference across FC paper and thesis;
% do not introduce synonyms (no "recovery-auth" where \axauth{} applies).
\newcommand{\axauth}{A1}          % Recovery authorisation  (formal: g)
\newcommand{\axrevault}{A2}       % Revault support         (formal: a)
\newcommand{\axfee}{A3}           % Fee-inclusion model     (formal: f)
\newcommand{\axwdbind}{A4}        % Withdrawal-destination binding (subcomp. of b)
\newcommand{\axrecbind}{A5}       % Recovery-destination binding   (subcomp. of b)
\newcommand{\axopsurf}{A6}        % Opcode surface          (new)
\newcommand{\axintro}{A7}         % Introspection primitive (new)

% Axis-with-formal-symbol aside (for proof-relevant passages).
\newcommand{\axauthf}{\axauth\,($g$)}
\newcommand{\axrevaultf}{\axrevault\,($a$)}
\newcommand{\axfeef}{\axfee\,($f$)}
\newcommand{\axbindf}{\axwdbind{}+\axrecbind\,($b$)}

% Attack-class labels --- canonical, see conference/AXES.md §4.
\newcommand{\classfee}{Z1}        % Fee-channel pinning
\newcommand{\classgrief}{Z2}      % Permissionless griefing
\newcommand{\classscale}{Z3}      % Per-UTXO recovery-cost scaling
\newcommand{\classhot}{Z4}        % Hot-key theft via destination substitution
\newcommand{\classcold}{Z5}       % Cold-key theft via unconstrained recovery
\newcommand{\classmode}{Z6}       % Mode/parameter bypass

% Draft-mode \todo — replaced or removed before submission
\newcommand{\todo}[2][]{\textbf{[TODO: #2]}}

% Measured-value placeholders. These render as [?] until the regtest
% harness produces the corresponding measurement (see
% vault-comparison/experiments/exp_watchtower_exhaustion.py Phase 6a).
% Replace each \measured{...} call with the numeric value before submission.
\newcommand{\measured}[1]{\textbf{[measured:#1]}}

% Batched-recovery constants — placeholders until Phase 6a runs on regtest
\newcommand{\overheadCCV}{54}
\newcommand{\perinputCCV}{68}
\newcommand{\overheadOPV}{154}
\newcommand{\perinputOPV}{92}

% ===================== Theorem environments =====================
% LNCS already defines theorem, definition, proposition, corollary, proof
% via \spnewtheorem in llncs.cls. No additional definitions needed.

% ===================== Paper metadata =====================
\title{The Cost of Custody}
\titlerunning{The Cost of Custody}

% Double-blind: authors anonymized until acceptance
\author{Anonymous Authors}
\authorrunning{Anonymous Authors}
\institute{Anonymous Institution\\
\email{anonymous@submission.invalid}}

\begin{document}

\maketitle

% ===================== Abstract =====================
\begin{abstract}
Bitcoin vaults encumber self-custodied funds with a time-delayed
spending path so that a compromised hot-key withdrawal can be
intercepted by a watchtower and redirected to cold storage. Four
script-extension proposals---\opctv{}, \opccv{}, \opvault{}, and
\catcsfs{}---are candidates for enabling this pattern, yet
practitioners choose among them without a peer-reviewed cross-design
security analysis.

We decompose covenant vault designs along four axes---fee model,
amount flexibility, recovery gating, and recovery binding---and show
that each axis choice deterministically opens a specific attack
surface. We instantiate all four designs on regtest under a uniform
lifecycle interface, run 15 experiments, and verify structural
properties using Alloy bounded model checking.

We report three quantitative results. First, prior hand-estimates
of \opvault{} transaction sizes are low by $45\%$, yielding a $36\%$
lifecycle cost premium over \opccv{}. Second, \opccv{}'s watchtower
splitting capacity reaches ${\sim}6{,}172$ splits per block---nearly
$2\times$ the widely-cited prior estimate---with single-block
exhaustion thresholds of $66\satvB$ for \opccv{} and $59\satvB$ for
\opvault{} at a $0.5$~BTC vault. Third, under batched recovery those
thresholds rise to $119\satvB$ and $159\satvB$, flipping the safety
ranking between the two designs---an operator-strategy dependence
prior work has not surfaced. We additionally prove a
griefing--safety incompatibility: no covenant vault simultaneously
achieves permissionless recovery and griefing resistance.

The cost of custody is therefore not a single number. It is a
function of design choice, threat model, fee regime, and watchtower
strategy, and no current proposal minimises it against every
adversary.
\end{abstract}

% ===================== Sections =====================
% §1 Introduction — Target 1.5pp

\section{Introduction}
\label{sec:intro}

Bitcoin's custody model places the entire burden of key security on the
asset holder: a compromised private key results in irreversible fund loss
with no protocol-level recourse. Covenant-based vaults address this by
interposing a time-delayed spending path between key compromise and fund
movement, so that an unauthorized hot-key spend can be intercepted by a
watchtower during a configurable delay window and redirected to cold
storage~\cite{SHMB20,MES16}. We focus on four Bitcoin Script
extension proposals with published BIPs, reference implementations,
and vault-specific threat discussion in the 2023--2026 activation
debate: \opctv{}~\cite{BIP119}, \opccv{}~\cite{BIP443}, \opvault{}~\cite{BIP345}
(withdrawn~\cite{OB25}, included for comparison), and the
\opcat{}+\opcsfs{} composition~\cite{BIP347,BIP348}. Other
covenant-capable proposals (ANYPREVOUT, TXHASH, TLUV) are beyond our
scope. Practitioners must choose among the four, yet no peer-reviewed
cross-design comparison exists.

The four proposals differ in how they introspect the spending
transaction, how they authorize recovery, and how they pay fees. These
differences are not cosmetic: each choice opens a distinct
vulnerability surface---anchor-based fee models admit descendant-chain
pinning (class \classfee{}), partial-withdrawal recovery admits a
per-UTXO recovery-cost tail whose ordering between \opccv{} and
\opvault{} inverts with watchtower batching (class \classscale{}),
keyless recovery admits permissionless griefing (class \classgrief{}),
destination-substitutable withdrawal admits hot-key theft
(class \classhot{}), and unbound recovery leaves admit
cold-key theft (class \classcold{}). No single design minimises
defender loss across all classes and all fee rates. The
activation debate~\cite{CTVCSFS25,Poinsot25,Towns25,OB25} has proceeded
without quantitative input on these tradeoffs; the literature contains
only qualitative matrices~\cite{covenants_info} and a single
back-of-envelope fee-pinning estimate~\cite{Har24} based on assumed
rather than measured transaction sizes.

This paper fills that gap. We build all four Bitcoin vault
implementations behind a uniform lifecycle interface and run 15
experiments on regtest measuring transaction sizes, attack costs,
and per-design functional differences. We additionally develop Alloy
bounded model checking specifications for structural properties (fund
conservation, recovery liveness, no-extraction), and derive
closed-form expressions for dust-threshold fee rates, pinning costs,
and batching amortisation. We implement a fifth design---a Simplicity
vault on Elements---as a cross-substrate reference point; because
Elements uses federated rather than proof-of-work consensus, the
Simplicity results are presented in Appendix~\ref{app:simplicity} and
are excluded from the formal statements below. Simplicity serves as
a constructive existence proof that the dominant corner of the
Bitcoin design space is occupiable on some substrate, even
while no Bitcoin proposal occupies it.

\paragraph{Contributions.}
\begin{enumerate}
  \item A \textbf{structural design-space decomposition}: covenant
    vault designs are characterised by seven canonical axes
    (\axauth{}, \axrevault{}, \axfee{}, \axwdbind{}, \axrecbind{},
    \axopsurf{}, \axintro{}; Section~\ref{sec:bg:axes}), with six
    derived attack classes \classfee{}--\classmode{}. Each class
    is a deterministic consequence of a specific axis-value
    (Sections~\ref{sec:imposs:class_fee}--\ref{sec:imposs:class_mode}),
    susceptibility and immunity both mechanical readings of the
    position table. No proposed Bitcoin extension occupies
    the dominant corner (Section~\ref{sec:imposs:design_space}).
  \item A \textbf{Griefing--Safety Incompatibility} proposition: no
    covenant vault simultaneously achieves permissionless recovery and
    griefing resistance (Section~\ref{sec:imposs:class_grief}).
  \item A \textbf{per-threat defender-loss characterisation}
    $\mathcal{L}(\mathrm{TM}, D, V, r, b)$ parameterised by threat
    model, design, vault value, fee rate, and defender batching
    strategy (Section~\ref{sec:imposs:sensitivity}).
  \item \textbf{Batched-defender measurements and a recovery-cost
    ordering-flip finding.} When the watchtower recovers many
    triggered vault UTXOs in one transaction instead of one at a
    time, its per-input cost drops sharply, raising the fee rate
    at which the per-UTXO recovery-cost tail of class \classscale{}
    saturates within a single block by $1.8\times$ for \opccv{} and
    $2.7\times$ for \opvault{}. The safer of the two designs under
    \classscale{} therefore depends on whether the watchtower
    batches: \opccv{} is safer without batching, \opvault{} is safer
    with it (Section~\ref{sec:empirical:batching}).
  \item \textbf{Axis-keyed deployment guidance} keyed to fee
    regime, batching strategy, and adversary profile, derived as
    counter-factual readings of the class rationality blocks
    (Section~\ref{sec:deployment}).
\end{enumerate}

Contributions are derived from the design-space decomposition of
Section~\ref{sec:bg:axes} and referenced uniformly throughout
Sections~\ref{sec:impossibility}--\ref{sec:deployment}. Every
attack-feasibility claim in this work is evaluated under the
Kerckhoffs threat model of Section~\ref{sec:bg:lifecycle}.

\paragraph{Scope.} Our formal results range over the four proposed
Bitcoin Script extensions only. We implement and discuss
Simplicity~\cite{OC17}, an alternative script primitive deployed on
the Elements federated sidechain~\cite{DPNS16}, as a cross-substrate
reference point (Appendix~\ref{app:simplicity}). Elements'
functionary-federated consensus differs fundamentally from Bitcoin's
proof-of-work, and the threats we analyse do not transfer identically
across that substrate boundary. Protocol-layer constructs such as Ark
and BitVM are also excluded: both operate above the opcode level and
require modelling interactive rounds outside our single-transaction
covenant framework.

\paragraph{Roadmap.}
Section~\ref{sec:background} summarises the four Script extensions,
the vault lifecycle, and the 11-element threat taxonomy.
Section~\ref{sec:impossibility} presents the structural results.
Section~\ref{sec:empirical} reports the regtest measurements.
Section~\ref{sec:deployment} derives deployment guidance from the
per-threat analysis. Section~\ref{sec:related} positions this work
against prior vault research and Section~\ref{sec:discussion}
discusses limitations and open problems.

\providecommand{\todo}[2][]{\textbf{[TODO: #2]}}

% §2 Background and Threat Model --- Target 2--2.5pp

\section{Background and Threat Model}
\label{sec:background}

\subsection{Covenants and the Four Introspection Primitives}
\label{sec:bg:covenants}

A \emph{covenant} is a spending condition that restricts not only who
can spend a UTXO but how the resulting transaction is
structured~\cite{MES16}. Implementing one requires \emph{transaction
introspection}: the ability of a script to examine fields of the
spending transaction. Every covenant proposal we study reduces to a
choice of \emph{introspection primitive}: the mechanism by which the
script observes transaction data. The four choices correspond to axis
\axintro{} in Section~\ref{sec:bg:axes} and determine each proposal's
downstream attack surface.

\paragraph{Digest equality (\opctv{}, BIP-119~\cite{BIP119}).} The
script carries a $32$-byte template hash. At spend time, the consensus
engine computes SHA-$256$ over the canonical concatenation
$\mathtt{nVersion}\,\|\,\mathtt{nLockTime}\,\|\,\mathtt{scriptSigsHash}\,\|\,\mathtt{inputCount}\,\|\,\mathtt{sequenceHash}\,\|\,\mathtt{outputCount}\,\|\,\mathtt{outputsHash}\,\|\,\mathtt{inputIndex}$
and checks equality against the committed hash. The primitive is
whole-template: any field divergence (including \texttt{outputsHash},
which itself commits to the \texttt{(amount,\,scriptPubKey)} pair of
every output) produces digest mismatch and rejection.

\paragraph{Merkle-membership with amount semantics (\opccv{},
BIP-443~\cite{Ing23}).} \opccv{} takes a mode byte selecting which
per-output check applies: $\{-1,\,0,\,1,\,2\}$ for self-input check,
output script check, output amount deduction, and combined checks
respectively. Recursion is supported via taptree preservation, which
enables partial unvault with automatic revault. Any other mode byte
triggers \texttt{OP\_SUCCESSx} and silently disables the covenant
check~(axis \axopsurf{}, threat TM8).

\paragraph{Dedicated opcodes (\opvault{}, BIP-345~\cite{BIP345}).}
\texttt{OP\_VAULT} (trigger) and \opvaultrecover{} (recovery) are
purpose-built opcodes with vault-specific semantics fixed at the
consensus layer. Amount preservation, automatic revault, and recovery
authorisation via a \texttt{recoveryauth} Schnorr key are enforced by
the opcode rather than by script composition. BIP-345 was withdrawn in
early 2025 in favour of BIP-443~\cite{OB25}; we retain it as the
purpose-built design point and a direct economic comparator for
\opccv{}.

\paragraph{Sighash-preimage reconstruction (\catcsfs{},
BIP-347+348~\cite{Poe21,BIP347,BIP348}).} \texttt{OP\_CAT} concatenates stack
items; \opcsfs{} verifies a Schnorr signature against a
stack-supplied message. The vault assembles the BIP-341 sighash
preimage from witness fields, then verifies the same signature under
the same public key against both the assembled preimage (via
\opcsfs{}) and the consensus-computed sighash (via
\texttt{OP\_CHECKSIG}). Soundness reduces to Schnorr's EUF-CMA under
the fixed-public-key model~\cite{NRS20}: a valid signature over two
distinct messages under the same key would imply a solution to the
discrete logarithm problem. Consequently the assembled preimage must
equal the consensus preimage, which binds the output set through
\texttt{hashOutputs}. The vault uses
\sighash{SINGLE}$|$\texttt{ACP} so that external fee-paying inputs
may be attached without breaking the binding. The recovery leaf,
however, is a bare \texttt{cold\_pk~OP\_CHECKSIG} with no
introspection, leaving the recovery destination
output-unconstrained~(axis \axrecbind{}, threat TM11).

\subsection{The Four Bitcoin Script Extension Proposals}
\label{sec:bg:designs}

Each proposal is a specific composition of the primitives above with
key-gating and fee-inclusion decisions. The dimensions that matter for
the analysis in Section~\ref{sec:impossibility} are enumerated as a
closed set of axes in Section~\ref{sec:bg:axes} below, and each
proposal's position on every axis is the canonical
Table~\ref{tab:design_space}.

\paragraph{\opctv{}.} The vault uses a two-step deterministic
transaction tree: a bare \opctv{} script guards the deposit, and the
unvault UTXO offers a hot withdrawal after a CSV delay or a
\opctv{}-committed cold sweep. An anchor output sized at $3{,}300$
sats supports CPFP fee bumping because the committed template cannot
anticipate future fee rates (axis \axfee{} =
\texttt{out-of-band-anchor}).

\paragraph{\opccv{}.} Per-output granularity enables partial unvault
with revault; the recovery leaf is \emph{keyless} (axis \axauth{} =
\texttt{keyless}), which creates a permissionless-griefing surface
(TM2) and is also the sole axis value that admits class \classgrief{}
as susceptible.

\paragraph{\opvault{}.} A three-key model separates the trigger key,
the \texttt{recoveryauth} key, and the recovery scriptPubKey.
\opvault{}'s value-preservation rule forces a separate fee-wallet
input on every trigger and recovery transaction, adding $80$--$90$\vB{}
of overhead but retaining dynamic fee inclusion (axis \axfee{} =
\texttt{dynamic}).

\paragraph{\catcsfs{}.} The dual-verification pattern (Section~\ref{sec:bg:covenants})
gives the strongest hot-key binding of the four designs (axis
\axwdbind{} = \texttt{signature-dual-bind}) but pairs it with the
weakest recovery binding (axis \axrecbind{} =
\texttt{plain-signature}); the cold-key
compromise surface (TM11) follows mechanically.

\subsection{The Design-Space Axes}
\label{sec:bg:axes}

The four proposals differ along seven axes. Three carry over from the
formal decomposition~\cite{SHMB20,Har24} and appear with Greek-letter
symbols ($f$, $a$, $g$, $b$) in proofs; the remaining four are
introduced in this work and are used as descriptive labels. Every
attack class in Section~\ref{sec:impossibility} derives from a
specific axis-value; every design's vulnerability profile is a lookup
against Table~\ref{tab:design_space}.

\begin{description}[leftmargin=1.5em,itemsep=1pt,parsep=0pt,topsep=2pt]
  \item[\axauthf{} --- Recovery authorisation.] Who may initiate the
    recovery path? Values: \texttt{keyless}, \texttt{recoveryauth-key},
    \texttt{hot-key}, \texttt{cold-key}. Controls class \classgrief{}
    (permissionless griefing).
  \item[\axrevaultf{} --- Revault support.] Does the covenant allow
    partial withdrawal? Values: \texttt{yes}, \texttt{no}. Controls
    class \classscale{} (per-UTXO recovery-cost scaling).
  \item[\axfeef{} --- Fee-inclusion model.] Where does the spending
    transaction's fee come from? Values: \texttt{dynamic} (fee in the
    transaction), \texttt{static-precommit} (fee baked at creation),
    \texttt{out-of-band-anchor} (CPFP via a separate anchor key).
    Controls class \classfee{} (fee-channel pinning).
  \item[\axwdbind{} --- Withdrawal-destination binding.] How is the
    withdrawal destination constrained? Values:
    \texttt{template-hash-digest}, \texttt{per-output-amount+script},
    \texttt{dedicated-opcode}, \texttt{signature-dual-bind}. Controls
    class \classhot{} (hot-key theft via destination substitution).
  \item[\axrecbind{} --- Recovery-destination binding.] Same question
    for the recovery path. Values: same as \axwdbind{}, plus
    \texttt{plain-signature} (no script-level constraint). Controls
    class \classcold{} (cold-key theft via unconstrained recovery).
    Kept separate from \axwdbind{} because \catcsfs{} occupies
    different positions on the two sub-axes --- strongest on
    \axwdbind{}, weakest on \axrecbind{}.
  \item[\axopsurf{} --- Opcode surface.] Is the covenant expressed by
    a single semantic opcode or a multi-mode polymorphic one with
    \texttt{OP\_SUCCESSx} fallthrough on undefined modes? Values:
    \texttt{single-mode}, \texttt{multi-mode-with-OP\_SUCCESS}.
    Controls class \classmode{} (mode/parameter bypass).
  \item[\axintro{} --- Introspection primitive.] How does the script
    read transaction data? Values correspond to the four primitives of
    Section~\ref{sec:bg:covenants}:
    \texttt{digest-equality},
    \texttt{merkle-membership-and-amount},
    \texttt{dedicated-opcode}, and
    \texttt{sighash-preimage-reconstruction}. Descriptive; governs
    which of the other axes are simultaneously realisable.
\end{description}

For the formal decomposition we retain $(f,\,a,\,g,\,b)$ as the axes
of the design-space lattice used in Propositions~1 and~2: $f =
\axfee$, $a = \axrevault$, $g = \axauth$, and $b$ is the joint
$(\axwdbind,\,\axrecbind)$ sublattice. Axes \axopsurf{} and \axintro{}
are additional dimensions orthogonal to the proofs; they are where the
developer-footgun (class \classmode{}) and the design-implementability
questions live.

\begin{table}[t]
\centering
\caption{The four Bitcoin covenant proposals, positioned on
the seven axes of Section~\ref{sec:bg:axes}. Every attack-class
susceptibility claim in Section~\ref{sec:impossibility} is a lookup
against this table. Simplicity (Appendix~D) occupies the strict-best
position on \axauth{}, \axrevault{}, \axwdbind{}, \axrecbind{}, and
\axintro{} simultaneously; no Bitcoin proposal does.}
\label{tab:design_space}
\footnotesize
\setlength{\tabcolsep}{2.2pt}
\renewcommand{\arraystretch}{1.12}
\begin{tabular}{@{}lllll@{}}
\toprule
Axis                                    & \opctv{}                       & \opccv{}                               & \opvault{}                    & \catcsfs{}                      \\
\midrule
\axauth{}  (recovery auth)              & hot-key                        & keyless                                & recoveryauth-key              & cold-key                        \\
\axrevault{} (revault)                  & no                             & yes                                    & yes                           & no                              \\
\axfee{}   (fee model)                  & out-of-band-anchor             & dynamic                                & dynamic                       & dynamic                         \\
\axwdbind{} (withdrawal binding)        & template-hash-digest           & per-output-amount+script               & dedicated-opcode              & signature-dual-bind             \\
\axrecbind{} (recovery binding)         & template-hash-digest           & per-output-amount+script               & dedicated-opcode              & plain-signature                 \\
\axopsurf{}  (opcode surface)           & single-mode                    & multi-mode-with-OP\_SUCCESS            & single-mode                   & single-mode                     \\
\axintro{}   (introspection primitive)  & digest-equality                & merkle-membership-and-amount           & dedicated-opcode              & sighash-preimage-reconstruction \\
\bottomrule
\end{tabular}
\renewcommand{\arraystretch}{1.0}
\end{table}

Two structural dependencies between axes constrain the realisable
combinations and are used implicitly throughout
Section~\ref{sec:impossibility}: (i) \axfee{} =
\texttt{out-of-band-anchor} implies \axrevault{} = \texttt{no}, because
the committed template cannot anticipate arbitrary future splits; and
(ii) \axintro{} = \texttt{sighash-preimage-reconstruction} is the only
known primitive that admits \axwdbind{} =
\texttt{signature-dual-bind}, because the binding reduces to Schnorr
message-equality under a single key.

\subsection{Vault Lifecycle, Threat Taxonomy, and Threat Model}
\label{sec:bg:lifecycle}

All four designs implement the Swambo et al.\ lifecycle~\cite{SHMB20}:
a deposit UTXO is triggered into an unvaulting UTXO subject to a
relative timelock (CSV); after the delay, the hot key completes the
withdrawal; at any point, the cold key (or, for \opccv{}, any party)
may invoke emergency recovery to a pre-committed address. Security
depends on an online watchtower detecting unauthorised triggers within
the CSV window.

\paragraph{Threat model (Kerckhoffs).} We adopt the standard
cryptographic threat model. The adversary knows the vault's design,
the covenant type, the script structure, and all public parameters
--- including the recovery address where applicable. Security rests on
keys and on the script's semantics, not on the secrecy of
construction. Every attack-feasibility claim in this work is evaluated
under this assumption. Concretely, we do not rely on taproot leaf
hiding, address non-disclosure, or closed-source construction as
security mechanisms.

\medskip We extend Swambo et al.\ with design-specific attack classes.
The 11-element threat taxonomy below serves as citation anchors for
the per-class analysis of Section~\ref{sec:impossibility} (classes
\classfee{}--\classmode{} subsume groups of TMs; see
Table~\ref{tab:design_space} and Appendix~B for the full class$\times$TM
map):

{\footnotesize
\begin{description}[leftmargin=1.3em,itemsep=0pt,parsep=0pt,topsep=2pt]
  \item[TM1 (fee-channel pinning).] Hot-key attacker saturates Bitcoin
    Core's $25$-descendant mempool cap with dust-rate transactions
    spending the \opctv{} unvault's anchor, blocking the watchtower's
    CPFP recovery. Class \classfee{}.
  \item[TM2 (recovery griefing).] An unauthorised party---anyone,
    under keyless recovery; the cold or auth-key holder, under
    key-gated---repeatedly invokes recovery, cycling the owner
    through re-deposit and re-trigger. Class \classgrief{}.
  \item[TM3 (trigger-key theft).] Hot-key attacker initiates an
    unauthorised unvault and races the watchtower during the CSV
    delay. Class \classhot{}.
  \item[TM4 (per-UTXO recovery-cost scaling).] Partial unvault with
    revault grows the vault's UTXO set over time; per-UTXO recovery
    cost scales linearly in the count and in fee rate. Under
    congestion, some UTXOs fall below the economic abandonment
    threshold. An attacker with the trigger key can accelerate this
    distribution; ordinary operation can produce it without one.
    Previously framed as ``watchtower exhaustion''; the reframe
    separates the structural cost tail from the multi-round attacker
    narrative (see Section~\ref{sec:empirical:batching} and
    Appendix~B).
    Class \classscale{}.
  \item[TM5 (trigger-xpub theft).] Capture of \opvault{}'s trigger
    xpub amplifies TM3 to every vault derived from that xpub.
  \item[TM6 (dual-key compromise).] Simultaneous theft of both the
    hot/trigger key and the cold/\texttt{recoveryauth} key defeats
    the split-key safety assumption.
  \item[TM7 (address reuse).] A second deposit to the same \opctv{}
    vault address produces a UTXO whose amount does not match the
    committed template hash, rendering it permanently unspendable. An
    operational-correctness failure rather than an adversarial attack
    (no adversary is required); included here as an axis-\axwdbind{}
    consequence observable under normal wallet use.
  \item[TM8 (\opccv{} mode bypass).] An undefined mode byte in the
    \opccv{} witness triggers \texttt{OP\_SUCCESSx}, silently
    disabling the covenant check. Specified behaviour under
    BIP-443's forward-compatibility design (axis \axopsurf{} =
    \texttt{multi-mode-with-OP\_SUCCESS}); the contribution is the
    systematic mode sweep and class-level framing, not the discovery.
    Class \classmode{}.
  \item[TM9 (hot-key redirection).] A \catcsfs{} hot-key attempt to
    redirect the withdrawal output is rejected by the dual
    \opcsfs{}+\texttt{OP\_CHECKSIG} binding; the outcome degrades to
    griefing only. Immunity derivation for class \classhot{}.
  \item[TM10 (witness tampering).] Modification of \catcsfs{}
    sighash-preimage witness fields is detected by the
    dual-verification check (stack-assembled preimage must agree with
    the consensus-computed sighash). Cryptographic verification,
    not an attack class.
  \item[TM11 (cold-key theft).] Capture of the cold/recovery key
    grants unrestricted fund redirection if the recovery path is
    output-unconstrained (\catcsfs{}); forced-recovery-only if
    output-bound (\opctv{}, \opvault{}). Class \classcold{}.
\end{description}
}

Three threats dominate the analysis: TM1 (\opctv{} anchor CPFP), TM2
(\opccv{} keyless recovery), and TM4 (per-UTXO recovery-cost scaling
on \opccv{}/\opvault{}). The per-design susceptibility matrix appears
in Section~\ref{sec:imposs:design_space}.

\providecommand{\todo}[2][]{\textbf{[TODO: #2]}}

% §3 Security Framework for Covenant Vault Designs
% Structure: per-class subsections (Z1–Z6), each deriving
% susceptibility and immunity from the axes of §2.3. Every proof
% preserved verbatim from the prior axis-structured draft.

\section{Security Framework for Covenant Vault Designs}
\label{sec:impossibility}

This section develops the security framework used throughout the
paper. Section~\ref{sec:imposs:framework} defines the defender-loss
metric that unifies theft, griefing, and dust-lockup outcomes.
Sections~\ref{sec:imposs:class_fee}--\ref{sec:imposs:class_mode}
develop six attack classes \classfee{}--\classmode{}; each class is a
mechanical consequence of a specific axis-value from
Table~\ref{tab:design_space}, with susceptibility and immunity
derived side-by-side rather than asserted.
Section~\ref{sec:imposs:design_space} identifies a structurally
dominant corner that no Bitcoin proposal occupies.
Section~\ref{sec:imposs:sensitivity} applies the framework
threat-by-threat, characterising how defender loss scales with fee
rate and recovery strategy.

% ================================================================
\subsection{Defender Loss and Adversary Advantage}
\label{sec:imposs:framework}

We parameterise each scenario by a threat model
$\mathrm{TM} \in \mathcal{T}$, design $D \in \mathcal{D}$, vault
value $V > 0$, fee rate $r \geq 0$, and defender strategy
$b \in \{0, 1\}$ ($b = 1$: defender batches recoveries; $b = 0$:
single-input). An adversary strategy $\mathcal{A}$ is a sequence of
on-chain transactions in pursuit of $\mathrm{TM}$.

\begin{definition}[Costs and outcomes]
\label{def:strategy}
For strategy $\mathcal{A}$ at fee rate $r$: $g(\mathcal{A}) \in [0,V]$
are extracted funds, $u(\mathcal{A}) \in [0,V]$ are dust-locked funds,
$c_A(\mathcal{A}, r)$ the adversary's fees, and $c_D(\mathcal{A}, r, b)$
the defender's fees under optimal response.
\end{definition}

\begin{definition}[Defender loss and adversary advantage]
\label{def:loss-adv}
Fix $(\mathrm{TM}, D, V, r, b)$ and attacker budget $B_A = V$
(rational-adversary assumption). The \emph{worst-case defender loss}
and \emph{adversary advantage} are
\[
  \mathcal{L} = \!\!\sup_{c_A(\mathcal{A}, r) \leq B_A}\!
    [\,g(\mathcal{A}) + u(\mathcal{A}) + c_D(\mathcal{A}, r, b)\,],
  \qquad
  \mathrm{Adv} = \sup_{\mathcal{A}} [\,g(\mathcal{A}) - c_A(\mathcal{A}, r)\,].
\]
$D_1 \prec_{\mathrm{TM}, V, r, b} D_2$ iff
$\mathcal{L}(\mathrm{TM}, D_1, V, r, b) < \mathcal{L}(\mathrm{TM}, D_2, V, r, b)$.
$\mathcal{L}$ unifies theft ($g>0$), griefing ($c_D > 0$), and
dust-lockup ($u > 0$); $\mathrm{Adv}$ is the attacker's rationality
signal and is $b$-invariant.
\end{definition}

Every covenant vault design is characterised by four independent
formal choices: a \emph{fee model}~$f$ (axis \axfee{}), an
\emph{amount-flexibility policy}~$a$ (axis \axrevault{}), a
\emph{recovery gating}~$g$ (axis \axauth{}), and a \emph{recovery
binding}~$b$ (jointly axes \axwdbind{} and \axrecbind{}). We write
$D = (f, a, g, b)$. Two spec-level dependencies link the axes: a fee
model that commits to a single output
(\sighash{SINGLE}|\texttt{ANYONECANPAY}) forces $a_{\text{atomic}}$;
conversely, $a_{\text{partial}}$ requires per-output introspection and
therefore $f \in \{f_{\text{wallet}}, f_{\text{inval}}\}$.
Table~\ref{tab:design_space} (Section~\ref{sec:bg:axes}) lists each
design's position on every axis; the six subsections that follow
derive the attack class enabled at each axis-value.

% ================================================================
\subsection{Class \classfee{}: Fee-Channel Pinning}
\label{sec:imposs:class_fee}

\paragraph{Class definition.} \classfee{} is the class enabled by
axis \axfee{} = \texttt{out-of-band-anchor}: fees attach via a CPFP
anchor output, whose own unconfirmed descendants block
defender-initiated replacement under standard-policy mempool rules.

\paragraph{The four fee models.} Four spec-level choices appear in
the proposals we study. \emph{CPFP anchor} ($f_{\text{anchor}}$,
\opctv{}): the unvault carries a small externally-spendable output
(${\approx}3{,}300$~sats) that lets any party bump the fee via
Child-Pays-For-Parent. \emph{Fee-wallet input} ($f_{\text{wallet}}$,
\opvault{}): value preservation forbids consuming vault value for
fees, so every trigger and recovery transaction attaches a separate
P2TR input from a fee wallet (adding 80--90\vB{}). \emph{In-value
deduction} ($f_{\text{inval}}$, \opccv{}): per-output introspection
allows fees to be deducted from the vault value itself. \emph{Sighash
ACP} ($f_{\text{acp}}$, \catcsfs{}): \sighash{SINGLE}|\texttt{ANYONECANPAY}
commits the signature to a single output and lets external
fee-paying inputs attach without breaking the covenant.

\paragraph{Susceptibility (implication D1).} \opctv{} is susceptible.
Under Bitcoin Core's default 25-descendant mempool relay policy, a
hot-key adversary broadcasts the unvault and attaches $25$ dust-rate
transactions spending the anchor. BIP-125 Rule 3 requires the
watchtower's CPFP replacement to pay for evicting every descendant
at a higher fee rate and higher absolute fee; against the attacker's
dust chain this is economically infeasible, blocking recovery and
forcing $\mathcal{L}(\mathrm{TM1}) = V$.

\paragraph{Immunity derivations.} \opccv{}, \opvault{}, and \catcsfs{}
have $\axfee{} \neq \texttt{out-of-band-anchor}$ and therefore carry
no pinnable anchor handle: dynamic fee inclusion keeps fee payment
inside the transaction the attacker would need to replace. The
attacker cannot pin a transaction whose existence they are trying to
delay. $\mathcal{L}(\mathrm{TM1}) = 0$ for all three.

\paragraph{Rationality (Appendix B.2).} \emph{Attacker capability:}
hot key (Kerckhoffs, §\ref{sec:bg:lifecycle}); capital for the
anchor-descendant chain ($\approx 14{,}300$~sats at $\rho = 3\satvB$).
\emph{Defender model:} online watchtower, Core-default mempool
policy. \emph{Preconditions:} $\mathcal{L}$ and $\mathrm{Adv}$ are
$r$-invariant. \emph{Counter-factual:} TRUC/BIP-431 collapses the
descendant limit to $1$, eliminating the pinning
surface~\cite{TRUC}; cluster mempool~\cite{ClusterMempool}
qualitatively changes the replacement calculus but is a scope-limit
rather than a within-scope mitigation. Full block:
Appendix~\ref{app:rationality}.

\paragraph{Empirical.} Section~\ref{sec:empirical:classes} reports
$C_{\text{pin}}$ and attack-composition with TM3.

% ================================================================
\subsection{Class \classgrief{}: Permissionless Griefing}
\label{sec:imposs:class_grief}

\paragraph{Class definition.} \classgrief{} is the class enabled by
axis \axauth{} = \texttt{keyless}: emergency recovery is invocable
without a signature, so any party --- including a pure griefer ---
can cycle the owner through re-deposit and re-trigger.

\paragraph{The two gating models.} Emergency recovery can be
\emph{keyless} ($g_{\text{keyless}}$, \opccv{}), invocable by any
party, or \emph{key-gated} ($g_{\text{key}}$,
\opctv{}\,/\,\opvault{}\,/\,\catcsfs{}), requiring a signature under a
specific cold or authorisation key. We formalise this choice via an
authorisation function
$\mathcal{A}\colon 2^{\mathcal{K}} \to \{0,1\}$ over the key set
$\mathcal{K}$: recovery is permissionless iff
$\mathcal{A}(\emptyset) = 1$.

\begin{proposition}[Griefing--safety incompatibility]
\label{thm:griefing_safety}
No covenant vault simultaneously achieves (i) permissionless
recovery ($\mathcal{A}(\emptyset) = 1$) and (ii) griefing
resistance (no unauthorised party can invoke recovery to force
$c_D > 0$).
\end{proposition}

\begin{proof}
Case analysis on $\mathcal{A}(\emptyset)$. If
$\mathcal{A}(\emptyset) = 1$, any third party meets the
authorisation predicate and can force recovery, requiring the
owner to re-deposit and re-trigger at cost $c_D > 0$; (ii) is
violated. If $\mathcal{A}(\emptyset) = 0$, unauthorised parties
cannot invoke recovery so (ii) holds, but (i) is violated by
definition. The cases exhaust
$\{0, 1\}$. \qed
\end{proof}

\paragraph{Susceptibility (implication D3).} \opccv{} is susceptible
and is the only design with \axauth{} = \texttt{keyless}. The
attacker pool is the entire population of Bitcoin nodes that can
observe the mempool; any such node can construct and broadcast a
valid recovery transaction consuming the unvaulting UTXO. The
Kerckhoffs assumption makes the recovery address public (cold
addresses are routinely disclosed for auditability or derived by the
same open-source library), so taproot leaf hiding does not provide
security here. $\mathcal{L}(\mathrm{TM2}) > 0$ for any fee rate.

\paragraph{Immunity derivations.} \opvault{} has \axauth{} =
\texttt{recoveryauth-key}: recovery demands a Schnorr signature
under a dedicated key, reducing the attacker pool to the key holder
and their compromises. \opctv{} has \axauth{} = \texttt{hot-key}
(recovery is the cold sweep, signed by the hot key with
\opctv{}-committed destination); the griefing surface collapses once
the hot key is safe. \catcsfs{} has \axauth{} = \texttt{cold-key}:
griefing requires cold-key compromise, which the attacker would
prefer to convert into class \classcold{} (Section~\ref{sec:imposs:class_cold}).

\paragraph{Rationality (Appendix B.2).} \emph{Attacker capability:}
zero keys; observation of mempool plus the public recovery address.
\emph{Defender model:} any; the attack is liveness-only and ends
only when the defender abandons or rotates the vault.
\emph{Preconditions:} $\mathcal{L}$ scales linearly in $r$;
$\gamma = \text{vsize}_{\text{atk}}/\text{vsize}_{\text{def}}$ is
$r$-invariant. \emph{Counter-factual:} upgrading \axauth{} to
\texttt{recoveryauth-key} (add a required signature on the recovery
leaf) eliminates the class; this is BIP-345's design choice.

\paragraph{Empirical.} Section~\ref{sec:empirical:classes} reports
$\gamma$ per design.

% ================================================================
\subsection{Class \classscale{}: Per-UTXO Recovery-Cost Scaling}
\label{sec:imposs:class_scale}

\paragraph{Class definition.} \classscale{} is the class enabled by
axis \axrevault{} = \texttt{yes}: partial withdrawal grows the
vault's UTXO count over time. Per-UTXO recovery cost scales linearly
in count and in fee rate; under congestion, some UTXOs fall below
the economic abandonment threshold. An attacker with the trigger
key can accelerate the distribution; ordinary operation can produce
it without one.

\paragraph{The two amount policies.} \emph{Atomic}
($a_{\text{atomic}}$, \opctv{}, \catcsfs{}): the full vault value
must exit the covenant in one transaction. \emph{Partial with
revault} ($a_{\text{partial}}$, \opccv{}, \opvault{}): per-output
introspection lets the spender produce a withdrawal portion and a
re-locked vault UTXO atomically, consuming one vault UTXO and
creating two outputs that between them preserve the original value.

\paragraph{Susceptibility (implication D2).} \opccv{} and \opvault{}
are susceptible. Repeated partial unvaults fragment the vault into
progressively smaller UTXOs. Each UTXO the defender must recover
costs $\text{vsize}_{\text{rec}} \cdot r$ in fees; when this exceeds
the UTXO's value, recovery is economically irrational and the UTXO
strands as dust. Above a fee-rate threshold derived in
Section~\ref{sec:imposs:sensitivity}, the defender loss
$\mathcal{L}(\mathrm{TM4})$ is strictly positive.

\paragraph{Immunity derivations.} \opctv{} and \catcsfs{} have
$a_{\text{atomic}}$ and cannot express partial withdrawal: the
template hash of \opctv{} and the sighash-preimage binding of
\catcsfs{} both commit to a fixed output set with no provision for
splitting. $\mathcal{L}(\mathrm{TM4}) = 0$ at every fee rate.

\paragraph{Reframing note.} Prior framings of this class as
``watchtower exhaustion'' cast it as a multi-round attacker-versus-
watchtower game. That framing does not survive a defender with full
recovery authority, who can full-sweep on first alert. The reframe
above identifies the class as a structural cost tail of partial
withdrawal --- attacker-accelerable but not attacker-required --- and
locates the defender's degree of freedom in recovery batching
(\ref{sec:empirical:batching}).

\paragraph{Rationality (Appendix B.2).} \emph{Attacker capability:}
trigger (hot) key (Kerckhoffs); capital for $N$ partial-withdrawal
triggers. \emph{Defender model:} watchtower with recovery authority;
batching flag $b$ governs per-UTXO cost scaling; minimal-disruption
policy (reactive per-UTXO recovery) maximises attacker leverage.
\emph{Preconditions:} rate-dependent; threshold
$r^\dagger_{\mathrm{TM4}}$ defined in
Section~\ref{sec:imposs:sensitivity}. \emph{Protocol assumptions:}
Bitcoin dust thresholds ($546$~sats P2WPKH, $330$~sats P2TR); RBF
rule 3 does not permit sibling-eviction, so the attacker's
partial-withdrawal chain is single-chain.
\emph{Counter-factual:} defender batching raises $r^\dagger$ by
$\mathord{\approx}\!1.8$--$2.7\times$
(Section~\ref{sec:empirical:batching}), absorbing the tail under
moderate congestion; full-sweep-on-first-alert policy collapses the
class to one round at the cost of operational disruption.

\paragraph{Empirical.} Section~\ref{sec:empirical:batching} reports
$(\alpha, \beta, r^\dagger)$ per design and batching flag, including
the \opccv{}/\opvault{} ordering flip that is the paper's novel
finding.

% ================================================================
\subsection{Class \classhot{}: Hot-Key Theft via Destination Substitution}
\label{sec:imposs:class_hot}

\paragraph{Class definition.} \classhot{} is the class enabled by
axis \axwdbind{} $\neq$ \texttt{signature-dual-bind}: the hot-key
holder, upon compromise, can direct the withdrawn funds to an
attacker-chosen destination, racing the watchtower during the CSV
delay. Signature-dual-binding on the withdrawal path uniquely
forecloses the redirection primitive at the cryptographic layer.

\paragraph{Susceptibility.} \opctv{} (\axwdbind{} =
\texttt{template-hash-digest}), \opccv{} (\axwdbind{} =
\texttt{per-output-amount+script}), and \opvault{} (\axwdbind{} =
\texttt{dedicated-opcode}) are all susceptible. In each case the
hot-key holder's signature authorises an unvault; the destination is
constrained by the \emph{script path spent} but the hot key can
select among alternative leaves or construct an unvault whose
destination the defender must recover from under the CSV race. The
outcome depends on output binding at unvault:
$\mathcal{L} = V \cdot \Pr_D[\text{attacker wins}] +
c_D^{\text{recover}} \cdot \Pr_D[\text{defender wins}]$.
$\Pr_D$ is dominated by $r$-independent factors (CSV length, network
topology), except that $\Pr_{\mathrm{OP\_CTV}}$ rises with TM1
applicability: class \classfee{} composes with \classhot{} under
\opctv{} to give $\mathcal{L} = V$ when fee pinning succeeds.

\paragraph{Immunity derivation: \catcsfs{}.} \catcsfs{} is immune by
\axwdbind{} = \texttt{signature-dual-bind}. Let $\sigma$ be the
Schnorr signature supplied in the witness, $m_{\text{preimage}}$ the
stack-assembled BIP-341 sighash preimage, and $m_{\text{consensus}}$
the sighash computed by consensus over the actual transaction.
\opcsfs{} checks
$\mathrm{Verify}(\mathrm{hot\_pk},\,m_{\text{preimage}},\,\sigma) = 1$;
\texttt{OP\_CHECKSIG} checks
$\mathrm{Verify}(\mathrm{hot\_pk},\,m_{\text{consensus}},\,\sigma) = 1$.
Both hold only if $m_{\text{preimage}} = m_{\text{consensus}}$: a
valid signature over two distinct messages under the same key would
violate Schnorr's EUF-CMA security~\cite{NRS20}, reducing to
discrete-log. The \texttt{hashOutputs} field is a component of both
preimages, so the output set is forced. Attempting to redirect yields
a rejected transaction, not a theft; the outcome degrades to
griefing-only (TM9). $\mathcal{L} = c_D$.

\paragraph{Rationality (Appendix B.2).} \emph{Attacker capability:}
hot key; mempool access; optional fee-pinning capital (if \classfee{}
composes). \emph{Defender model:} online watchtower able to race
during CSV; race bias toward attacker under TM1 composition.
\emph{Preconditions:} CSV length, network topology; $r$-independent
except via TM1.
\emph{Counter-factual:} \axwdbind{} = \texttt{signature-dual-bind}
(as in \catcsfs{}) eliminates the class by cryptographic construction;
any covenant able to express dual-verification can port the
primitive.

\paragraph{Empirical.} Section~\ref{sec:empirical:classes} reports
$\mathcal{L}$ per design and the composition with \classfee{}.

% ================================================================
\subsection{Class \classcold{}: Cold-Key Theft via Unconstrained Recovery}
\label{sec:imposs:class_cold}

\paragraph{Class definition.} \classcold{} is the class enabled by
axis \axrecbind{} = \texttt{plain-signature}: the recovery path is a
bare \texttt{cold\_pk OP\_CHECKSIG} with no script-level
output constraint. Cold-key compromise produces unrestricted fund
redirection rather than forced recovery to the pre-committed
destination.

\paragraph{The two recovery-binding models.} \emph{Output-bound}
($b_{\text{bound}}$, \opctv{}, \opvault{}, \opccv{}, Simplicity): the
recovery output scriptPubKey is fixed at vault creation (via a
\opctv{} template, a \opvault{}/\opccv{}-committed destination, or a
Simplicity output-hash jet). \emph{Unbound}
($b_{\text{unbound}}$, \catcsfs{} recovery leaf): the recovery leaf
is a bare \texttt{cold\_pk OP\_CHECKSIG} with \sighash{DEFAULT},
placing no constraint on outputs.

\paragraph{Susceptibility (implication D4).} \catcsfs{} is
susceptible. A compromised cold key produces a valid signature over
any sighash, which \texttt{OP\_CHECKSIG} accepts; there is no
script-level check on \texttt{hashOutputs}.
$\mathcal{L}(\mathrm{TM11}) = V$; the attacker chooses the
destination freely. $r$-invariant.

\paragraph{Immunity derivations.} \opctv{} and \opvault{} have
\axrecbind{} \,$\in$\, \{\texttt{template-hash-digest},
\texttt{dedicated-opcode}\}: the recovery script commits to a
specific recovery destination, so a compromised cold key can force
recovery but cannot redirect funds.
$\mathcal{L}(\mathrm{TM11}) = 0$. \opccv{} is excluded by
\axauth{} = \texttt{keyless} (no cold key to compromise); its
\axrecbind{} = \texttt{per-output-amount+script} would likewise
prevent redirection.

\paragraph{Note (axis \axwdbind{} vs.\ \axrecbind{} asymmetry).}
\catcsfs{}'s strongest axis position is \axwdbind{} =
\texttt{signature-dual-bind} (uniquely immune to \classhot{}) paired
with its weakest axis position \axrecbind{} = \texttt{plain-signature}
(uniquely susceptible to \classcold{}). Splitting $b$ into the two
sub-axes \axwdbind{} and \axrecbind{} (§\ref{sec:bg:axes}) turns
this asymmetry from a remark into a mechanical reading of the
position table. \cite{Poe21} sketches a \catcsfs{}-based bound
recovery leaf that would move the design to \axrecbind{} =
\texttt{signature-dual-bind}, eliminating \classcold{}.

\paragraph{Rationality (Appendix B.2).} \emph{Attacker capability:}
cold key (Kerckhoffs --- the attack assumes compromise, not
inference). \emph{Defender model:} any; the attack is instant with no
CSV window. \emph{Counter-factual:} \axrecbind{} =
\texttt{signature-dual-bind} on the recovery leaf eliminates the
class; alternatively, \axauth{} = \texttt{keyless} prevents the
cold-key compromise surface but is Pareto-dominated by Proposition~1.

\paragraph{Empirical.} Section~\ref{sec:empirical:classes}.

% ================================================================
\subsection{Class \classmode{}: Mode/Parameter Bypass}
\label{sec:imposs:class_mode}

\paragraph{Class definition.} \classmode{} is the class enabled by
axis \axopsurf{} = \texttt{multi-mode-with-OP\_SUCCESS}: a
polymorphic covenant opcode takes a stack-provided mode selector,
and the taproot \texttt{OP\_SUCCESSx} upgrade mechanism promotes
undefined mode values to unconditional script success. A developer
or construction bug that selects an undefined mode silently disables
the covenant check.

\paragraph{Susceptibility.} \opccv{} is susceptible. BIP-443 defines
four valid modes $\{-1, 0, 1, 2\}$; any other value falls through to
\texttt{OP\_SUCCESSx}, which under BIP-341/342 taproot rules
terminates script execution with success regardless of stack state.
Five undefined modes $\{3, 4, 7, 128, 255\}$ were measured on
regtest; each produces full covenant bypass at $110$\vB{} per spend.
$\mathcal{L}(\mathrm{TM8}) = V$.

\paragraph{Immunity derivations.} \opctv{}, \opvault{}, and \catcsfs{}
have \axopsurf{} = \texttt{single-mode}: each exposes a single
semantic opcode with no mode selector, so the \texttt{OP\_SUCCESSx}
fallthrough surface does not arise.

\paragraph{Framing.} The class is \emph{specified behaviour} under
BIP-443's forward-compatibility design (itself inherited from
BIP-341/342), not a newly-discovered vulnerability. The contribution
here is the systematic mode sweep and the class-level framing that
ties the surface to axis \axopsurf{}, not the discovery of
\texttt{OP\_SUCCESSx} semantics. In practice, the pymatt reference
library mitigates the risk by providing named constants
(\texttt{CCV\_FLAG\_*}); the residual risk is in independent
implementations that hand-write the mode byte.

\paragraph{Rationality (Appendix B.2).} \emph{Attacker capability:}
any observer of a vault whose script contains an undefined mode
byte; no keys required. \emph{Defender model:} developer discipline
around mode-byte selection; library-provided named constants.
\emph{Counter-factual:} \axopsurf{} = \texttt{single-mode} (one
opcode per primitive, as in \opctv{}, \opvault{}, \catcsfs{})
eliminates the class. A bitmask-gated version of \opccv{} (reject
undefined modes at script-parse time) would also eliminate the
class without requiring a soft-fork revision, but would forfeit the
forward-compatibility benefit.

\paragraph{Empirical.} Appendix~\ref{app:threat_matrix} documents the
five-mode sweep.

% ================================================================
\subsection{The Dominant Corner and Current Proposals}
\label{sec:imposs:design_space}

Combining the axes, each design occupies one cell of the
lattice. Restricting to the proof-relevant formal sublattice
$(f, a, g, b)$, define the \emph{dominant corner}
$D^\star = (f_{\neq\text{anchor}}, a_{\text{atomic}},
g_{\text{key}}, b_{\text{bound}})$. By implications D1--D4, any
design at $D^\star$ has $\mathcal{L} = 0$ simultaneously against
\{TM1, TM2, TM4, TM11\} --- equivalently, against classes
\{\classfee{}, \classgrief{}, \classscale{}, \classcold{}\}.
Table~\ref{tab:pareto-summary} locates each proposed BIP in the
sublattice with Hamming distance to $D^\star$; the full axis
positions are Table~\ref{tab:design_space} in
Section~\ref{sec:bg:axes}.

\begin{table}[t]
\centering
\caption{The four Bitcoin Script extensions in the $(f, a, g, b)$
sublattice of Table~\ref{tab:design_space}. Distance measures Hamming
separation from the dominant corner $D^\star$ on the proof-relevant
sublattice; the full A1--A7 position of each design is available in
Section~\ref{sec:bg:axes}.}
\label{tab:pareto-summary}
\small
\begin{tabular}{@{}lccccc@{}}
\toprule
Design     & $f$                      & $a$                  & $g$                     & $b$                    & dist.\ \\
\midrule
\opctv{}        & $f_{\text{anchor}}$      & $a_{\text{atomic}}$  & $g_{\text{key}}$        & $b_{\text{bound}}$     & 1 \\
\opccv{}        & $f_{\text{inval}}$       & $a_{\text{partial}}$ & $g_{\text{keyless}}$    & $b_{\text{bound}}$     & 2 \\
\opvault{} & $f_{\text{wallet}}$      & $a_{\text{partial}}$ & $g_{\text{key}}$        & $b_{\text{bound}}$     & 1 \\
\catcsfs{}   & $f_{\text{acp}}$         & $a_{\text{atomic}}$  & $g_{\text{key}}$        & $b_{\text{unbound}}$   & 1 \\
$D^\star$  & $\neq f_{\text{anchor}}$ & $a_{\text{atomic}}$  & $g_{\text{key}}$        & $b_{\text{bound}}$     & --- \\
\bottomrule
\end{tabular}
\end{table}

No proposed Bitcoin BIP occupies $D^\star$. \opctv{} differs
in~$f$ and inherits \classfee{} by D1; \opccv{} differs in~$a$
and~$g$ and inherits \classscale{} and \classgrief{} by D2 and D3;
\opvault{} differs in~$a$ and inherits \classscale{} by D2;
\catcsfs{} differs in~$b$ and inherits \classcold{} by D4. A
hypothetical specification at $D^\star$ would have $\mathcal{L} = 0$
against all four of \{\classfee{}, \classgrief{}, \classscale{},
\classcold{}\} simultaneously; the corner is currently open on
Bitcoin. Two present proposals are one axis away: \opctv{}
(by replacing its anchor-based fee model, which TRUC
adoption~\cite{TRUC} approximates) and \catcsfs{} (by replacing its
bare cold-key recovery leaf with an output-bound path, sketched
in~\cite{Poe21}). Simplicity on Elements (Appendix~D) occupies the
dominant corner --- demonstrating the corner is \emph{occupiable},
just not by any current Bitcoin proposal.
\label{thm:fee_inversion}\label{cor:open_corner}

\paragraph{Localisation for future proposals.}
The framework yields a lookup procedure: read a proposal's
specification, extract its axis tuple
$(\axauth, \axrevault, \axfee, \axwdbind, \axrecbind, \axopsurf, \axintro)$,
and consult the class derivations
(\ref{sec:imposs:class_fee}--\ref{sec:imposs:class_mode}) to read
off its vulnerability profile. No further measurement is needed to
characterise which of \{\classfee{}, \classgrief{}, \classscale{},
\classhot{}, \classcold{}, \classmode{}\} a newly-proposed BIP
admits by construction. What the framework does \emph{not}
determine---numerical thresholds like $r^\dagger_{\mathrm{TM4}}(V,
D, b)$---depends on user parameters ($V$, watchtower block budget)
and measured transaction sizes, and is the subject of
Section~\ref{sec:imposs:sensitivity} and the empirical measurements
of Section~\ref{sec:empirical}.

% ================================================================
\subsection{Per-Threat Loss Characterisation}
\label{sec:imposs:sensitivity}

We now apply the framework threat by threat, deriving closed-form
expressions for $\mathcal{L}$ as a function of fee rate and
defender strategy. Each TM is the scope-anchor of the class noted in
parentheses; see Appendix~\ref{app:rationality} for the full
seven-field rationality blocks.

\paragraph{TM1 (fee-channel pinning, class \classfee{})~\cite{Har24}.}\label{sec:sens:tm1}
Applies to \opctv{} alone by D1.
$\mathcal{L}(\mathrm{TM1}, \mathrm{OP\_CTV}, V, r) = V$;
$\mathrm{Adv} = V - C_{\text{pin}}$ with
$C_{\text{pin}} = 25(\rho+1)\,\text{vsize}_{\text{desc}} + \alpha_{\text{anchor}}
\approx 14{,}300$~sats ($\rho = 3\satvB$, anchor $= 3{,}300$~sats).
Both $\mathcal{L}$ and $\mathrm{Adv}$ are $r$-invariant;
$\mathcal{L} = 0$ for $D \neq \mathrm{OP\_CTV}$.

\paragraph{TM2 (recovery griefing, class \classgrief{}).}
Applies to all designs, with attacker pools set by D3. Over $k$
rounds,
\begin{equation} \label{eq:gamma}
  \mathcal{L}(\mathrm{TM2}, D, V, r) = k \cdot \text{vsize}_{\text{def}}(D) \cdot r,
  \qquad
  \gamma(D) = \frac{\text{vsize}_{\text{atk}}(D)}{\text{vsize}_{\text{def}}(D)}.
\end{equation}
$\mathcal{L}$ scales linearly in $r$; $\gamma$ is $r$-invariant.
Ordering by $\gamma$: \opccv{} ($0.27$, keyless, severe) $<$ \catcsfs{}
($0.36$) $<$ \opvault{} ($0.55$) $<$ \opctv{} ($0.71$).

\paragraph{TM3 (trigger-key theft, class \classhot{}).}
Applies to all designs. The hot-key adversary unvaults and races
the watchtower during the CSV delay; outcome depends on output
binding at unvault:
$\mathcal{L} = V \cdot \Pr_D[\text{attacker wins}] +
c_D^{\text{recover}} \cdot \Pr_D[\text{defender wins}]$. $\Pr_D$ is
dominated by $r$-independent factors (CSV length, network
topology), except that $\Pr_{\mathrm{OP\_CTV}}$ rises with TM1
applicability. Ordering: \catcsfs{} ($\mathcal{L} = c_D$,
dual-verification binding on the hot leaf) $<$ \opvault{} $\approx$
\opccv{} ($\mathcal{L} \in [c_D, V]$) $<$ \opctv{} ($\mathcal{L} = V$ when
TM1 succeeds).

\paragraph{TM4 (per-UTXO recovery-cost scaling, class \classscale{}).}\label{sec:sens:tm4}
Applies to \opccv{} and \opvault{} by D2. Let $\text{vsize}_{\text{rec}}(D)$
be the single-input recovery vsize. Under batched recovery, BIP-345
§``Batching'' and BIP-443 default aggregation permit the defender
to recover $N$ UTXOs in one transaction of vsize
$\alpha(D) + \beta(D) \cdot N$. The effective per-UTXO recovery
cost is
\begin{equation}
  \text{vsize}_{\text{rec}}(D, b, N) =
  \begin{cases}
    \text{vsize}_{\text{rec}}(D)                             & \text{if } b = 0, \\
    \alpha(D)/N + \beta(D)                                   & \text{if } b = 1.
  \end{cases}
\end{equation}
A split becomes dust when its value falls below this effective
cost; the number of splits needed to drive every UTXO below dust
is $N_{\text{dust}}(D, r, b) = \lfloor V /
(\text{vsize}_{\text{rec}}(D, b, N_{\text{dust}}) \cdot r) \rfloor$.
Let $S(D)$ denote splits per block. The single-block-feasibility
threshold
\begin{equation} \label{eq:r-dagger}
  r^\dagger_{\mathrm{TM4}}(V, D, b)
  \;=\; \frac{V}{\text{vsize}_{\text{rec}}(D, b, \cdot) \cdot S(D)}
\end{equation}
is the fee rate at which $N_{\text{dust}}/S = 1$; lower
$r^\dagger$ means the structural cost tail saturates at lower fees,
hence less safe. Batching strictly raises $r^\dagger$ whenever
$\beta(D) < \text{vsize}_{\text{rec}}(D)$, always true in practice.
Measured $(\alpha, \beta)$ values, the resulting thresholds, and
the \opccv{}/\opvault{} ordering (which flips between $b = 0$ and
$b = 1$) appear in Section~\ref{sec:empirical:batching}. The
multi-round attacker narrative of prior framings is one
threat-model instantiation among several; the class-level statement
is that the cost tail exists as a consequence of $a_{\text{partial}}$
independently of adversary behaviour.

\paragraph{TM11 (cold-key compromise, class \classcold{}).}\label{sec:sens:tm11}
Applies to key-gated designs by D4.
$\mathcal{L}(\mathrm{TM11}) = V$ under $b_{\text{unbound}}$
(\catcsfs{}: bare \texttt{cold\_pk OP\_CHECKSIG} with
\sighash{DEFAULT});
$\mathcal{L}(\mathrm{TM11}) = 0$ under $b_{\text{bound}}$ (\opctv{},
\opvault{}). \opccv{} is excluded by $g_{\text{keyless}}$ (no cold key
to compromise). $r$-invariant in both cases.

\paragraph{Protocol-level vs.\ user-parameter-level.}
The \emph{existence} of the threshold $r^\dagger_{\mathrm{TM4}}$
is a protocol-level consequence of D2. Its numerical value depends
on $V$ (user-chosen vault value), $N_{\text{budget}}$ (defender
block budget), and $\text{vsize}_{\text{rec}}$ fixed by BIP
specification; implementation variation is below 5\vB{} in our
measurements. The implications D1--D4 thus establish
protocol-level structural properties; reported numerical values
scale linearly with user parameters.

\begin{table}[t]
\centering
\caption{Per-class worst-case defender loss $\mathcal{L}$ across
the four Bitcoin covenant proposals. Values are for
$V = 0.5$~BTC. Simplicity is treated separately (Appendix~D).
Every zero-loss cell is annotated with the structural reason,
keyed to the footnote below.
\textbf{Legend:} $V$ = full theft; $0$ = design immune by the
structural property noted;
$\gamma = \text{vsize}_{\text{atk}} / \text{vsize}_{\text{def}}$;
$[c_D, V]$ = race-dependent bounded interval;
$r^\dagger$ = single-block-feasibility fee rate (\satvB{}).
Fee-rate dependence: inv = invariant, $\propto r$ = linear,
thr = threshold, race = race-dependent.}
\label{tab:pareto-matrix}
\renewcommand{\arraystretch}{1.1}
\begin{tabular}{@{}lccccc@{}}
\toprule
Threat (class)                                     & $r$-dep.    & \opctv{}                    & \opccv{}                  & \opvault{}       & \catcsfs{}               \\
\midrule
TM1 pinning (\classfee{})                          & inv         & $V$                    & $0^{\mathrm{a}}$     & $0^{\mathrm{b}}$ & $0^{\mathrm{c}}$       \\
TM2 griefing (\classgrief{})                       & $\propto r$ & $\gamma{=}0.71$        & $\gamma{=}0.27$      & $\gamma{=}0.55$  & $\gamma{=}0.36^{*}$    \\
TM3 trig-theft (\classhot{})                       & race        & $[c_D, V]^{\dagger}$   & $[c_D, V]$           & $[c_D, V]$       & $c_D$                  \\
TM4 cost scaling (\classscale{})                   & thr         & $0^{\mathrm{d}}$       & $r^\dagger{=}66$     & $r^\dagger{=}59$ & $0^{\mathrm{e}}$       \\
TM11 cold key (\classcold{})                       & inv         & $0$                    & $0^{\mathrm{f}}$     & $0$              & $V$                    \\
\bottomrule
\end{tabular}
\par\vspace{3pt}
\footnotesize
\textbf{Zero-loss cells explained:}
$^{\mathrm{a,b,c}}$ no anchor output, hence no descendant-chain
pinning surface; $^{\mathrm{d,e}}$ no partial-unvault-with-revault
primitive (atomic withdrawal only), so no split-exhaustion surface;
$^{\mathrm{f}}$ keyless recovery path, no cold key to compromise.
$^{*}$\catcsfs{}'s TM2 attack requires the cold key; if that key is
compromised the attacker prefers TM11 (full theft). $^{\dagger}$TM3
composes with TM1 to achieve $V$ when fee pinning succeeds.
\end{table}

No row of Table~\ref{tab:pareto-matrix} is dominated in every
column, and TM4's fee-rate-dependence interacts with TM1's
invariance to make aggregate rankings unstable. We return to the
resulting deployment surface in Section~\ref{sec:deployment}.

% §4 Empirical Validation --- Target 3pp

\section{Empirical Validation}
\label{sec:empirical}

\subsection{Methodology}
\label{sec:emp:method}

A Docker-based regtest harness runs each vault design under a
uniform lifecycle interface; experiments invoke the vault
operations (deposit, trigger, withdraw, recover) without knowledge
of the underlying covenant mechanism and dispatch on per-design
capability flags rather than names. Adapters bind to the
required node variant: Bitcoin Inquisition (\opctv{}, \catcsfs{}),
Merkleize \opccv{}~\cite{BIP443}, jamesob's \opvault{} branch, and
Elements under Simplex (Simplicity, Appendix~\ref{app:simplicity}).
The primary metric is virtual size (vsize): structurally
deterministic, identical on regtest and mainnet. Fees are
projections $\text{fee} = \text{vsize} \times r$ since regtest has
no fee market. Temporal recovery races are not captured (blocks mine
on demand). Fifteen experiments span lifecycle costs, batching,
revault chains, address reuse, class \classfee{} (fee-channel
pinning), class \classscale{} (per-UTXO recovery-cost scaling), and
the 11-element threat taxonomy; measurements are reproducible via the
anonymised artifact (DOI withheld for double-blind review).

\subsection{Lifecycle Transaction Sizing}
\label{sec:emp:sizes}

Table~\ref{tab:vsizes} reports the measured vsize for every
transaction in the deposit--trigger--withdraw lifecycle. Each value
is structurally deterministic (range~=~0 across repeated runs). The
rightmost column tags the axis value whose consequences each row
measures, per Section~\ref{sec:bg:axes}.

\begin{table}[t]
\centering
\caption{Lifecycle vsize measurements (\vB{}) for the four
Bitcoin covenant proposals. Lifecycle total~=~deposit +
trigger + withdraw. Simplicity on Elements is presented in
Appendix~\ref{app:simplicity}. The \emph{axis-value} column gives
the primary axis (Table~\ref{tab:design_space}) whose position the
row's size measures as a consequence.}
\label{tab:vsizes}
\begin{tabular}{@{}lrrrrcl@{}}
\toprule
Design      & deposit & trigger & withdraw & recover & \textbf{total} & axis-value                                          \\
\midrule
\opctv{}    & 122     &  94     & 152      & 133     & \textbf{368}   & \axfee{}=anchor, \axwdbind{}=digest                 \\
\opccv{}    & 300     & 154     & 111      & 122     & \textbf{565}   & \axrevault{}=yes, \axfee{}=dynamic                  \\
\opvault{}  & 154     & 292     & 121      & 246     & \textbf{567}   & \axrevault{}=yes, \axfee{}=dynamic (wallet)         \\
\catcsfs{}  & 122     & 221     & 210      & 125     & \textbf{553}   & \axwdbind{}=dual-bind, \axfee{}=dynamic (ACP)       \\
\bottomrule
\end{tabular}
\end{table}

\opctv{} has the cheapest lifecycle (368\vB{}); its trigger is a bare
template-hash check with no signature (94\vB{}, smallest of any
design). \opccv{} and \opvault{} total nearly identically (565 vs.\
567\vB{}) with opposite distributions: \opccv{} has a cheap trigger
(154\vB{}) and expensive deposit (300\vB{}, two contract outputs);
\opvault{} has an expensive trigger (292\vB{}, fee-wallet input) and
moderate deposit. \catcsfs{}'s dual \opcsfs{}+CHECKSIG pattern falls
between (553\vB{}), costing trigger 221\vB{} and withdraw 210\vB{}.

\paragraph{\opvault{} vsize correction.}
Prior informal estimates of \opvault{} \texttt{start\_withdrawal}
(${\sim}200$\vB{}) and recovery (${\sim}170$\vB{}) are both under by
$\approx 45\%$ (measured: 292\vB{}, 246\vB{}). BIP-345's
value-preservation rule forces a P2TR fee-wallet input
($80$--$90$\vB{}) on every non-deposit transaction, yielding a 36\%
lifecycle cost premium over \opccv{}---an economic factor in the BIP-345
withdrawal~\cite{OB25}.

\subsection{Class Measurements: \classfee{} and \classscale{}}
\label{sec:emp:attacks}
\label{sec:empirical:classes}

We reproduce the two primary economic attack classes on regtest. For
class \classfee{} (fee-channel pinning, TM1,
Section~\ref{sec:imposs:class_fee}), we construct the 25-chain
descendant attack against a \opctv{} unvault's 3{,}300-sat anchor and
confirm $C_{\text{pin}} \approx 14{,}300$~sats---$0.03\%$ of a
$0.5$~BTC vault, rational at any fee regime. For class \classscale{}
(per-UTXO recovery-cost scaling, TM4,
Section~\ref{sec:imposs:class_scale}), substituting the measured
$\text{vsize}_{\text{rec}}$ of Table~\ref{tab:vsizes} into
Eq.~\ref{eq:r-dagger} yields unbatched thresholds
$r^\dagger_{\mathrm{OP\_CCV}} \approx 66$\satvB{} and
$r^\dagger_{\opvault{}} \approx 59$\satvB{} at $V = 0.5$~BTC. Above
$r^\dagger$ the structural cost tail saturates within a single block;
below it the defender can absorb the per-UTXO liability before
dust-economics kick in. \opctv{} and \catcsfs{} are immune to
\classscale{} by \axrevault{} = \texttt{no} (atomic withdrawal only).

\subsection{Batched Recovery and the Ordering Flip}
\label{sec:empirical:batching}

This is the paper's empirically novel claim: under class \classscale{},
the safety ranking between \opccv{} and \opvault{} inverts based on
whether the watchtower batches recovery. The multi-round attacker
scenario of prior framings is one threat-model instantiation
among several (Appendix~\ref{app:rationality}); the ordering-flip
result holds independently of attacker behaviour and is a statement
about \emph{defender implementation choice}.

BIP-345 §``Batching'' and BIP-443 default aggregation permit the
defender to recover $N$ triggered UTXOs in a single transaction. We
swept $N \in \{2, 5, 10, 25, 50\}$ on our regtest harness for each
batching-capable design; linear regression yields
$\text{vsize}_{\text{rec}}(D, N) = \alpha(D) + \beta(D) \cdot N$:
\begin{itemize}
  \item \opccv{}: $\alpha_{\mathrm{OP\_CCV}} = \overheadCCV$~\vB{},
    $\beta_{\mathrm{OP\_CCV}} = \perinputCCV$~\vB{}.
  \item \opvault{}: $\alpha_{\opvault{}} = \overheadOPV$~\vB{},
    $\beta_{\opvault{}} = \perinputOPV$~\vB{}.
\end{itemize}

Substituting $\beta(D)$ for $\text{vsize}_{\text{rec}}(D)$ in
Eq.~\ref{eq:r-dagger}, the batched thresholds at $V = 0.5$~BTC rise
to 119\satvB{} (\opccv{}) and 159\satvB{} (\opvault{}), factors of
$1.8\times$ and $2.7\times$ over unbatched. The ordering flips: \opccv{}
is safer unbatched ($66 > 59$) but \opvault{} is safer batched
($159 > 119$). \opvault{}'s larger unbatched recovery transaction
amortises disproportionately; design choice under \classscale{} thus
depends on the defender's batching strategy, and an ordering stated
without reference to $b$ is incomplete.

\paragraph{Asymmetric benefit.}
BIP-345 authorised recovery permits batching across vaults with
distinct \texttt{<recovery-sPK-hash>} values; unauthorised recovery
is capped at two outputs and thus restricted to
same-\texttt{recovery-sPK} batches. \opccv{}'s default aggregation has no
taptree-equivalence requirement and is strictly more permissive. The
attacker's partial-unvault-with-revault sequence is inherently
serial---each revault consumes the previous output---so the
attacker does not benefit from batching, and the measured savings
accrue only to the defender.

\subsection{Simplicity on Elements}
\label{sec:simplicity-perspective}

Our Simplicity vault on Elements supports per-output binding on
trigger and recovery, atomic or bounded partial withdrawal, cold-key
recovery with an optional keyless post-timeout sweep, and MAX\_FEE-capped
in-value fee deduction. In the atomic configuration with
cold-key-only recovery, it occupies the dominant corner of
Section~\ref{sec:imposs:design_space} (strict-best on \axauth{},
\axrevault{}, \axwdbind{}, \axrecbind{}, and \axintro{}
simultaneously). Empirical measurements (lifecycle 865\vB{} for the
atomic configuration; threat coverage; Coq-verified semantics) and
the bounded-partial and keyless-post-timeout constructions are
reported in Appendix~\ref{app:simplicity}. Its role in the argument
is to demonstrate that the dominant corner is occupiable on some
substrate, even while no Bitcoin proposal currently occupies it.

\providecommand{\todo}[2][]{\textbf{[TODO: #2]}}

% §5 Deployment Guidance --- Target 1.5pp

\section{Deployment Guidance}
\label{sec:deployment}

The per-class analysis of Sections~\ref{sec:imposs:class_fee}--\ref{sec:imposs:class_mode}
and the empirical thresholds of Section~\ref{sec:empirical} yield a
deployment matrix keyed on four parameters: (i) fee regime, (ii)
whether the watchtower batches recovery, (iii) whether the deployment
prioritises hot-key theft resistance or key-loss resilience, and (iv)
whether the deployment tolerates a federated sidechain for the
dominant-corner properties. All recommendations are derived as
counter-factual readings of the class rationality blocks
(Appendix~\ref{app:rationality}); none are asserted outside that
structure. All recommendations assume the Kerckhoffs threat model of
Section~\ref{sec:bg:lifecycle}: the attacker is taken to know the
deployment's axis values.

\begin{table}[t]
\centering
\caption{Deployment recommendations keyed on fee regime, watchtower
batching capability, and the deployment's priority (hot-key theft
resistance vs.\ key-loss resilience). Each recommendation is a
counter-factual reading of the class rationality blocks. ``Defer
to corner'' = none of the Bitcoin proposals is
class-free in this cell; operator must pick by risk weighting, or
accept federated trust to deploy on Simplicity
(Appendix~\ref{app:simplicity}).}
\label{tab:deployment}
\footnotesize
\setlength{\tabcolsep}{2.5pt}
\renewcommand{\arraystretch}{1.15}
\begin{tabular}{@{}lllll@{}}
\toprule
Fee regime                           & Batching        & Priority              & Recommended    & Rationale (class counter-factual)                                                   \\
\midrule
Low ($\leq 10\satvB{}$)              & either          & theft-resistance      & \catcsfs{}     & \classhot{} immune by \axwdbind{} dual-bind; \classscale{} N/A at low $r$           \\
Low ($\leq 10\satvB{}$)              & either          & key-loss resilience   & \opccv{}       & \classgrief{} low-cost at low $r$; key-loss covered by \axauth{} = keyless          \\
Mid ($10$--$100\satvB{}$)            & batched         & either                & \opvault{}     & \classscale{} mitigated by batching (\classscale{} §counter-factual)                \\
Mid ($10$--$100\satvB{}$)            & unbatched       & theft-resistance      & \catcsfs{}     & As above; unbatched defenders cannot safely absorb \classscale{} above $r^\dagger$  \\
Mid ($10$--$100\satvB{}$)            & unbatched       & key-loss resilience   & Defer to corner & No axis-A1 = keyless design without \classscale{} susceptibility among Bitcoin proposals \\
High ($\geq 100\satvB{}$)            & batched         & either                & \opvault{}     & Batched $r^\dagger \approx 159\satvB{}$; no anchor to pin                           \\
High ($\geq 100\satvB{}$)            & unbatched       & either                & \opctv{}       & \classfee{} cost $r$-invariant; \classscale{} N/A by \axrevault{} = no              \\
High, post-TRUC~\cite{TRUC}          & either          & either                & \opctv{}       & \classfee{} eliminated at relay-policy level; no other class active                 \\
Any                                  & any             & all classes zero      & Simplicity     & Appendix~\ref{app:simplicity}; accepts federated trust                              \\
\bottomrule
\end{tabular}
\end{table}

\paragraph{Reading the table.} Each row is a deployment scenario; the
Recommended column cites the design whose class exposures zero-out
most efficiently under the row's conditions; the Rationale column
names the specific counter-factual from the corresponding class
rationality block. Class \classfee{} (TM1) is eliminated \emph{at the
relay-policy level} by TRUC/BIP-431~\cite{TRUC}: when a \opctv{}
unvault uses a TRUC transaction type, the descendant limit collapses
to $1$ and the pinning surface vanishes. This is the single most
impactful deployment decision an \opctv{} operator can make.

\paragraph{Watchtower reserve sizing.}
For a $V = 0.5$~BTC vault at a fee rate of $100\satvB{}$ defending
against roughly one thousand unbatched \opccv{} splits, the recovery
budget is about $12.2$M~sats (${\sim}24\%$ of the vault value); the
same scenario under batching costs about $6.8$M~sats
(${\sim}14\%$), a direct consequence of the $1.8\times$
improvement in the single-block-feasibility threshold. Operators
should size reserves against the 95th-percentile fee rate of the
target deployment window and test the emergency-recovery path
quarterly.

\paragraph{Recovery-key storage policy.}
Proposition~\ref{thm:griefing_safety} forces a binary choice on axis
\axauth{}. Operators who prioritise fund safety under key loss
(\axauth{} = \texttt{keyless}, \opccv{}) accept griefing exposure
(class \classgrief{}) and must budget for the recovery-griefing
scenario (Appendix~\ref{app:rationality}, \classgrief{}). Operators
who prioritise griefing resistance (\axauth{} $\neq$
\texttt{keyless}, \opctv{}/\opvault{}/\catcsfs{}) reintroduce
key-loss risk and must use hardware-grade cold-key protection ---
particularly acute for \catcsfs{}, whose recovery leaf is
output-unconstrained (\axrecbind{} = \texttt{plain-signature},
class \classcold{}).

\paragraph{What the table does not say.}
Numerical thresholds ($r^\dagger$, reserve percentages) are
\emph{derived}, not \emph{axiomatic}. They depend on $V$, on the
watchtower's block budget, and on measured transaction sizes
(Table~\ref{tab:vsizes}). An operator with a different vault value
or a different blocks-per-recovery budget must resubstitute into
Eq.~\ref{eq:r-dagger} rather than treat this table's thresholds as
constants. The class derivations
(Sections~\ref{sec:imposs:class_fee}--\ref{sec:imposs:class_mode})
are the operator-independent content; the table is the
operator-parameterised specialisation for the reference
$V = 0.5$~BTC case.

\providecommand{\todo}[2][]{\textbf{[TODO: #2]}}

% §6 Related Work

\section{Related Work}
\label{sec:related}

\paragraph{Vault custody protocols.}
Swambo et al.~\cite{SHMB20} formalised the vault lifecycle---deposit,
unvault, withdraw, recover---and the watchtower monitoring assumption
that underpins all covenant vault designs, establishing the threat
vocabulary we adopt and extend to the 11-element taxonomy. Their
companion work~\cite{SBMB20} characterised three implementation
approaches (deleted-key pre-signed transactions, COSHV, \opctv{}) but
did not compare opcode-level designs empirically. Swambo's doctoral
thesis~\cite{Swa23} extended this to the Ajolote autonomous custody
system but did not include \opctv{}, \opccv{}, \opvault{}, or
\catcsfs{}. Swambo and Poinsot's risk framework for the Revault
protocol~\cite{SP21} developed an operational risk taxonomy for
pre-signed-transaction vaults; our per-threat analysis is in a
similar spirit but targets opcode-level covenants rather than
deleted-key designs. Möser, Eyal, and Gün Sirer~\cite{MES16}
introduced covenants as a vault-custody mechanism, predating all
four BIPs. O'Beirne's technical report on vaults~\cite{OBeirne23}
informed the design of BIP-345, which was subsequently
withdrawn~\cite{OB25}. Kirstein et al.~\cite{KLGZ21} formally
verified a single regenerating vault design in Dafny; our work
differs in comparing five designs under a uniform methodology
rather than deeply verifying one.

\paragraph{Covenant proposals and analysis.}
Each covenant proposal has its own design documentation and
community analysis~\cite{BIP119,BIP443,BIP345,BIP347,BIP348}, but no
peer-reviewed work compares them empirically. Harding~\cite{Har24}
published a back-of-envelope estimate of \opvault{} single-block
recovery-cost saturation (historically framed as ``watchtower
exhaustion'') at ${\sim}3{,}000$ splits per block based on
assumed---not measured---transaction sizes; we confirm $3{,}427$ for
\opvault{} and measure $6{,}172$ for \opccv{}, revealing \opccv{}'s
roughly $2\times$ worse unbatched exposure. Poelstra~\cite{Poe21}
demonstrated \opcat{}-based introspection with dynamic destination
and recursive covenants; our \catcsfs{} vault implements a
simplified variant (fixed destination, no recursion) that isolates
the dual-verification security properties. Rijndael~\cite{Rij24}
built a Schnorr-G-trick vault without \opcsfs{}, demonstrating that
\opcat{} alone enables covenant enforcement.

\paragraph{Formal models of Bitcoin contracts.}
The Bitcoin smart-contract formalisation lineage---Atzei et al.'s
formal transaction model~\cite{ABLZ18}, BitML and its
extensions~\cite{BZ18,ABL19,BLM22}, BitML's covenant
extension~\cite{BLZ20}, and the Ethereum-side SoK
methodology~\cite{ABC17}---provides the process-algebraic substrate
against which our Alloy state-machine models stand as a
complementary abstraction level, operating over the $(f, a, g, b)$
design tuple. dgpv~\cite{dgpv24} applied Alloy to Rijndael's
\opcat{} vault; we extend that methodology to all four BIP-specified
designs plus Simplicity under a shared base model, adding a
threat-model layer and a cross-covenant module.
O'Connor~\cite{OC17} provided Coq-verified denotational semantics
for Simplicity; our Simplicity model operates at the state-machine
level, complementing that work with application-level verification.

\paragraph{Layer-two security and rational adversaries.}
Gudgeon et al.'s SoK~\cite{GMRMG20} systematises layer-two
protocols; watchtower-based delegated monitoring has been
formalised in TEE~\cite{LGSM19}, on-chain
arbitration~\cite{McCorry19}, and UC-style state-channel
constructions~\cite{DFH18,DEFM19,ABEFHMM21,Aumayr24}. Our class \classscale{}
analysis shares the always-online monitoring assumption but reframes
the question from ``how does an attacker exhaust the watchtower?''
to ``what is the structural per-UTXO recovery-cost tail of partial
withdrawal, and how does defender batching interact with it?'' On the
rational-adversary side, miner-bribery analyses against timelocked
Bitcoin primitives~\cite{TYME21,WSZN23,NKW21} formalise the
adversary class that covenant vaults defend against during the
unvault delay. Riard~\cite{Ria23} demonstrated replacement cycling
against Lightning HTLCs, a mempool-manipulation class that may
apply to covenant recovery transactions and is out of scope here.

\paragraph{Positioning.}
The qualitative covenants.info matrix~\cite{covenants_info}
compares proposals across feature dimensions without transaction
sizes, attack costs, or fee analysis. No prior work combines
empirical measurement, formal verification, and
fee-environment-dependent security analysis across multiple
covenant vault designs. Our contribution is the integration of
these three threads under a uniform adapter-based methodology
that enables both Proposition~\ref{thm:griefing_safety} and the
design-space framework of Section~\ref{sec:imposs:design_space}.

\providecommand{\todo}[2][]{\textbf{[TODO: #2]}}

% §7 Discussion and Limitations — Target 1pp

\section{Discussion and Limitations}
\label{sec:discussion}

\subsection{Limitations}
\label{sec:disc:limitations}

\paragraph{Kerckhoffs scope.}
Our threat model (Section~\ref{sec:bg:lifecycle}) assumes the
attacker knows the vault's axis values, including public parameters
such as the recovery address. We do not rely on hiding attacks
(taproot leaf concealment, address non-disclosure, closed-source
construction) as security mechanisms: a vault whose security depends
on the attacker's ignorance of its construction is not
cryptographically secure by Kerckhoffs's criterion. Any result or
recommendation in this paper that is read outside this threat model
should be reinterpreted accordingly.

\paragraph{Regtest vs.\ mainnet.}
All measurements are taken on Bitcoin regtest with a single mempool.
The class \classfee{} rationality block
(Appendix~\ref{app:rationality}) scopes the fee-pinning feasibility
claim: in practice it depends on a network-layer propagation race
between the attacker's 25 descendants and the watchtower's CPFP
replacement. Regtest cannot model this race. Structural vsize
measurements are unaffected (they are identical on regtest and
mainnet), but the \emph{feasibility} of a race-sensitive attack
depends on deployment-environment parameters beyond our harness.
Mainnet-equivalent measurement on Inquisition signet is future work.

\paragraph{Bounded model checking.}
Alloy's finite-scope analysis (41 assertions across five models)
provides refutation for broad claims but does not constitute a full
proof. Mechanised proofs (Coq, Lean) of
Proposition~\ref{thm:griefing_safety} and of the feature-vulnerability
implications D1--D4 (Section~\ref{sec:impossibility}) remain open.

\paragraph{Unmerged BIPs.}
BIP-119, BIP-345, BIP-443, BIP-347, and BIP-348 are all under
review, and specifications may evolve. We tracked specifications via
pinned git hashes (see anonymised artifact); readers should verify
against the current BIP text.

\paragraph{Class \classhot{} (TM3) race dynamics are network-layer,
not Script-layer.}
Our design-space decomposition does not model the race between the
attacker's unvault-plus-withdrawal sequence and the watchtower's
recovery broadcast. Formalising $\Pr_D[\text{attacker wins race}]$
requires a separate probability space over block-arrival times,
mempool propagation delays, and detection latency---parameters that
depend on the deployment environment rather than on BIP script
semantics. Existing state-channel watchtower formalisations treat
confirmation as atomic or assume bounded-delay synchrony; none
provides off-the-shelf machinery for a pinning-composed race against
a Bitcoin mempool. We therefore treat class \classhot{} (TM3) as
uniformly applicable to all four Script extensions at the script
layer, with \catcsfs{}'s dual-binding on the hot leaf
(\axwdbind{} = \texttt{signature-dual-bind}) as the one design-level
mitigation (reducing the worst-case outcome from theft to griefing).

\paragraph{Batching operational complexity.}
The measured $r^\dagger$ gains of
Section~\ref{sec:empirical:batching} characterise the best-case
defender strategy. Real deployments face buffering (CSV-window
deadlines per pending UTXO), single-input-poisons-batch failure
modes, and, for BIP-345 authorised recovery, an online
\texttt{recoveryauth} key that enlarges the watchtower's own attack
surface. These are upper bounds on achievable safety margin, not
refutations of the thresholds.

\subsection{Open Problems}
\label{sec:disc:open}

\begin{enumerate}
  \item Mainnet-equivalent measurement on Inquisition signet.
  \item Mechanised proofs (Coq/Lean) of
    Proposition~\ref{thm:griefing_safety} and of the
    feature-vulnerability implications D1--D4
    (Section~\ref{sec:impossibility}).
  \item Impact of TRUC/BIP-431 relay policy on class \classfee{}
    (fee-channel pinning). Under TRUC the 25-descendant limit
    collapses to 1, eliminating the pinning surface at the
    relay-policy level (see Section~\ref{sec:imposs:class_fee}
    counter-factual and Appendix~\ref{app:rationality}).
  \item Recursive covenant design space beyond \opccv{}'s taptree
    preservation.
  \item Cross-chain generalisation: Ethereum account abstraction as
    an L1 analogue of the design-space decomposition.
\end{enumerate}

\providecommand{\todo}[2][]{\textbf{[TODO: #2]}}

% §8 Conclusion — Target 0.5pp

\section{Conclusion}
\label{sec:conclusion}

We presented the first empirical comparison of four Bitcoin
covenant proposals for vault custody---\opctv{}, \opccv{},
\opvault{}, \catcsfs{}---plus a Simplicity vault on Elements
(Appendix~\ref{app:simplicity}) as a cross-substrate reference that
occupies the dominant corner. Two structural results constrain the
design space. Proposition~\ref{thm:griefing_safety} shows that no
covenant vault simultaneously achieves permissionless recovery and
griefing resistance. The design-space decomposition of
Section~\ref{sec:bg:axes} shows that vault designs factor along
seven canonical axes, with six derived attack classes
\classfee{}--\classmode{} each a deterministic consequence of a
specific axis-value (Sections~\ref{sec:imposs:class_fee}--\ref{sec:imposs:class_mode}).
As a structural consequence, no proposed Bitcoin extension
sits at the corner of the space where defender loss is
simultaneously zero against all six classes; each falls short on
at least one axis with a corresponding class surface. The
\axauth{}--\axintro{} axes provide a lookup procedure for localising
future BIPs without additional measurement.

The cost of custody is therefore not a single number. It depends on
the covenant's axis tuple, the class the operator prioritises, the
prevailing fee regime, and the watchtower's operational strategy.
Under two designs equal at the script layer, the relative safety can
still invert with the defender's strategy: safety rankings among
current covenants are not operator-free. No proposal minimises the
cost of custody against every class; the seven-axis decomposition
makes the tradeoffs precise and identifies a structurally-safer
corner $D^\star$ that the Bitcoin BIP corpus has not yet occupied.
Open directions include mechanised proofs of the class derivations,
the impact of TRUC/BIP-431 and cluster-mempool on class \classfee{},
temporal verification of class \classhot{} races, and cross-chain
generalisation of the decomposition. The framework, all measurements,
formal models, and both original vault implementations are available
at the anonymised artifact (URL withheld for double-blind review).

\providecommand{\todo}[2][]{\textbf{[TODO: #2]}}


% ===================== Bibliography =====================
\bibliographystyle{splncs04}
\bibliography{references}

% ===================== Appendix =====================
\appendix
% Appendix A: Alloy Formal Models

\section{Alloy Formal Models}
\label{app:alloy}

The bounded model checking referenced in Section~\ref{sec:discussion}
uses Alloy 6 with a three-layer module hierarchy: a Bitcoin base
layer (\texttt{btc\_base}), a vault abstraction layer
(\texttt{vault\_base}), and five concrete covenant models extending
the abstraction with design-specific transition guards. Two
cross-cutting modules (\texttt{threat\_model},
\texttt{cross\_covenant}) are opened by the concrete models. In
total the corpus contains 42 assertions; all but one pass under the
configured scopes.

\subsection{Module Contents}

\paragraph{\texttt{btc\_base.als}.}
Signatures for \texttt{UTXO}, \texttt{Key}, \texttt{Transaction},
and \texttt{Address}. UTXOs carry an amount, owner key, and
spent-status flag; transactions consume inputs and produce outputs.

\paragraph{\texttt{vault\_base.als}.}
Abstract vault state machine with four states (\texttt{VAULTED},
\texttt{UNVAULTING}, \texttt{WITHDRAWN}, \texttt{RECOVERED}),
transitions (trigger, withdraw, recover), and a \texttt{VaultFamily}
relation grouping UTXOs of a single vault instance. The CSV delay
is an integer field on each vault UTXO. Invariants defined at this
level (fund conservation, recovery liveness) propagate to all
concrete models.

\paragraph{Concrete vault models.}
Five modules extend \texttt{vault\_base} with covenant-specific
transition guards: \texttt{ctv\_vault.als} models the template-hash
commitment; \texttt{ccv\_vault.als} models the mode enumeration and
\texttt{OP\_SUCCESS} behaviour for undefined modes;
\texttt{opvault\_vault.als} models the three-key separation and
authorised recovery; \texttt{cat\_csfs\_vault.als} models the
dual-verification binding and the unconstrained recovery leaf;
\texttt{simplicity\_vault.als} models the
\texttt{jet::outputs\_hash()} check on both trigger and recovery
paths.

\paragraph{\texttt{threat\_model.als}.}
Attacker profiles parameterised by which keys the attacker
controls; each property is checked under each profile to identify
the minimum capability needed for each attack class.

\paragraph{\texttt{cross\_covenant.als}.}
Models concurrent vaults from different covenant types sharing the
same UTXO set; checks cross-vault injection and fee-wallet
contention between concurrent \opvault{} operations.

\subsection{Asserted Invariants}
\label{app:alloy:invariants}

The main shared invariants are:

\begin{description}[leftmargin=1.4em,itemsep=1pt,parsep=0pt,topsep=2pt]
  \item[Fund conservation.] Total vault value is preserved across
    lifecycle transitions except via recovery or withdrawal.
  \item[Destination binding.] The withdraw path produces only the
    pre-committed withdrawal output (\opctv{}, \opvault{}, Simplicity);
    for \opccv{} the withdraw path is \opctv{}-committed; \catcsfs{}'s hot
    leaf is bound via dual verification.
  \item[Recovery authorisation.] Recovery can only be invoked by a
    holder of the design's authorisation set
    ($\emptyset$ for \opccv{}; cold key for \opctv{}, \catcsfs{}, Simplicity;
    \texttt{recoveryauth} key for \opvault{}).
  \item[Recovery liveness.] From \texttt{VAULTED} or
    \texttt{UNVAULTING}, some transition to \texttt{RECOVERED}
    exists under the authorisation set.
  \item[No cross-leaf extraction.] Spending the trigger leaf
    cannot produce outputs that would satisfy the recovery leaf
    and vice versa.
\end{description}

\subsection{Per-Model Verification Results}

\paragraph{\opctv{}.} Six assertions all pass, including
no-extraction under each single-key compromise profile, CSV
enforcement, and the no-state-proliferation property. The
address-reuse scenario finds an instance where a second deposit
produces an unspendable UTXO, confirming TM7 as structural.

\paragraph{\opccv{}.} Eight assertions pass, including a
mode-confusion containment check (undefined modes trigger
\texttt{OP\_SUCCESS} but only within the affected leaf; the
vault's other leaves remain secure). The griefing scenario
demonstrates that any party may invoke recovery without keys
(TM2). The splitting scenario exhibits the revault chain that
produces the per-UTXO recovery-cost tail of class \classscale{}
(TM4).

\paragraph{\opvault{}.} Eight assertions pass, including a
dual-key-no-theft property: joint compromise of the trigger and
\texttt{recoveryauth} keys yields liveness denial but not theft,
because the recovery address is pre-committed in the Taproot
tweak. A key-loss scenario exhibits permanent inability to recover
if the \texttt{recoveryauth} key is destroyed.

\paragraph{\catcsfs{}.} Five of six assertions pass. The
\texttt{catcsfsNoExtraction\_ColdKeyOnly} assertion \emph{fails}:
Alloy exhibits a trace in which the cold-key holder invokes
recovery with an attacker-chosen output address, extracting funds.
This is the structural confirmation of TM11 and the only expected
assertion failure in the corpus. Destination-lock and CSV
assertions pass.

\paragraph{Simplicity.} Six assertions pass: no-extraction under
any single-key compromise, output binding enforcement on both
trigger and recovery, CSV enforcement, and no state proliferation.
A recovery-constrained scenario confirms that cold-key compromise
can only recover to the pre-committed address, unlike \catcsfs{}.

\subsection{Scope and Limitations}

The Alloy models operate at the state-machine level, abstracting
cryptographic primitives (signatures are modeled as key-possession
checks; hash functions are assumed collision-resistant),
serialisation, and temporal dynamics. Network propagation, mining,
and mempool policy are out of scope. Verification confirms
\emph{structural} reachability and authorisation properties but
not cryptographic unforgeability or temporal race outcomes; the
regtest experiments (Section~\ref{sec:empirical}) complement the
formal layer with measured transaction sizes and empirical attack
costs.

\providecommand{\todo}[2][]{\textbf{[TODO: #2]}}

% Appendix B: Full Threat Matrix + Rationality and Scope Blocks
%
% B.1 is the per-design qualitative severity matrix (expanded form of
% Table 5 in Section §3); B.2 is the class rationality and scope
% matrix that populates the seven-field rationality block format of
% AXES.md §7 for each of the six attack classes Z1-Z6. The FC paper
% compresses the rationality blocks to matrix form; the thesis App D
% carries the full free-text form.

\section{Complete Threat Model Matrix}
\label{app:threat_matrix}
\label{app:threat-matrix}

\subsection{Per-Design Severity (TM1--TM11)}
\label{app:threat_matrix:severity}

Table~\ref{tab:pareto-matrix} in Section~\ref{sec:imposs:design_space}
reports worst-case defender loss $\mathcal{L}$ across five dominant
threats (TM1, TM2, TM3, TM4, TM11). Table~\ref{tab:threat-matrix-full}
below lists the remaining threats (TM5--TM10) alongside the dominant
ones for completeness, using a qualitative severity scale
(\textbf{Safe} / \textbf{N/A} / \textbf{Low} / \textbf{Moderate} /
\textbf{High} / \textbf{Severe}) and citing measured vsize where
applicable. Simplicity is excluded from this matrix because
Elements' substrate alters the threat semantics
(Section~\ref{sec:simplicity-perspective}).

\begin{table}[h]
\centering
\caption{Per-design severity for the full 11-element threat
taxonomy. Severity labels: \textbf{N/A} = structurally
inapplicable; \textbf{Low} = requires difficult preconditions;
\textbf{Moderate} = feasible under specific conditions;
\textbf{High} / \textbf{Severe} = significant risk under realistic
assumptions; \textbf{Safe} = correct-by-design; \textbf{Impossible}
= prevented by a cryptographic or structural invariant.}
\label{tab:threat-matrix-full}
\footnotesize
\setlength{\tabcolsep}{3pt}
\begin{tabular}{@{}cp{2.6cm}p{2.6cm}p{2.6cm}p{2.6cm}@{}}
\toprule
TM & \opctv{} & \opccv{} & \opvault{} & \catcsfs{} \\
\midrule
1 & \textbf{High.} 25-descendant chain blocks CPFP (${\approx}14{,}300$~sats).
  & \textbf{N/A} (no anchor).
  & \textbf{N/A} (fee-wallet input).
  & \textbf{N/A} (external fee input). \\[4pt]

2 & \textbf{Moderate.} Hot-key triggers; defender sweeps via cold leaf.
  & \textbf{Severe.} Keyless; $\gamma{=}0.27$.
  & \textbf{Low.} Auth key required; $\gamma{=}0.55$.
  & \textbf{Low.} Cold key required; subsumed by TM11. \\[4pt]

3 & \textbf{High.} Hot+fee $\Rightarrow$ theft; hot alone $\Rightarrow$ DoS.
  & \textbf{High.} Redirect; watchtower races (122\vB).
  & \textbf{Moderate.} Redirect; authorised recovery races (246\vB).
  & \textbf{Low.} Griefing only (dual-verification binding). \\[4pt]

4 & \textbf{N/A} (no partial unvault).
  & \textbf{Severe.} $\sim$6{,}172 splits/block; $r^\dagger \approx 66\satvB$.
  & \textbf{Severe.} $\sim$3{,}427 splits/block; $r^\dagger \approx 59\satvB$.
  & \textbf{N/A} (no partial unvault). \\[4pt]

5 & \textbf{N/A} (no xpub derivation).
  & \textbf{N/A} (no xpub derivation).
  & \textbf{Moderate.} xpub expansion leaks all children.
  & \textbf{N/A} (no xpub derivation). \\[4pt]

6 & \textbf{High.} Hot + fee $\Rightarrow$ theft.
  & \textbf{N/A} (single key; keyless recovery).
  & \textbf{High.} Trigger + auth $\Rightarrow$ persistent DoS.
  & \textbf{High.} Cold key alone $\Rightarrow$ theft (via TM11). \\[4pt]

7 & \textbf{High.} Stuck funds on address reuse.
  & \textbf{Safe.}
  & \textbf{Safe.}
  & \textbf{Safe.} \\[4pt]

8 & \textbf{N/A} (no mode byte).
  & \textbf{Developer footgun.} Undefined mode $\rightarrow$ \texttt{OP\_SUCCESS}.
  & \textbf{N/A} (no mode byte).
  & \textbf{N/A} (no mode byte). \\[4pt]

9 & \textbf{N/A} (no \opcsfs{} path).
  & \textbf{N/A} (no \opcsfs{} path).
  & \textbf{N/A} (no \opcsfs{} path).
  & \textbf{Impossible.} Dual verification rejects redirection. \\[4pt]

10 & \textbf{N/A} (no \opcsfs{} witness).
   & \textbf{N/A} (no \opcsfs{} witness).
   & \textbf{N/A} (no \opcsfs{} witness).
   & \textbf{Impossible.} SHA-256 binding. \\[4pt]

11 & \textbf{N/A} (\opctv{}-committed recovery address).
   & \textbf{N/A} (keyless recovery; no cold key).
   & \textbf{N/A} (Taproot-tweak-constrained recovery).
   & \textbf{High.} Bare \texttt{cold\_pk OP\_CHECKSIG} $\Rightarrow$ theft. \\
\bottomrule
\end{tabular}
\end{table}

\paragraph{Simplicity (cross-substrate reference).}
Under honest-federation assumptions on Elements, TM1 does not
transfer (no public mempool / 25-descendant cap). Our Simplicity
vault is atomic (\axrevault{} = \texttt{no}), so TM4 has no surface;
a partial-withdrawal variant would inherit \classscale{} per the D2
implication (Section~\ref{sec:imposs:class_scale}). TM11 outcome
matches \opctv{} and \opvault{} (output-bound recovery via
\texttt{jet::outputs\_hash()}). TM5, TM8, TM9, TM10 do not apply:
Simplicity has no xpub derivation, no mode byte, no \opcsfs{} path,
and no \opcsfs{} witness. TM7 is safe (per-vault plan). TM6 reduces
to TM11 since the hot path is also output-bound. Full design-space
position and lifecycle measurements:
Appendix~\ref{app:simplicity}.

\subsection{Rationality and Scope Blocks}
\label{app:rationality}

Each attack class \classfee{}--\classmode{} is accompanied by a
seven-field rationality block per the schema of Sections
\ref{sec:imposs:class_fee}--\ref{sec:imposs:class_mode}. For the FC
paper we compress the blocks into Table~\ref{tab:rationality} with
one cell per (class, field) intersection. The thesis Appendix D
carries the uncompressed free-text form.

Every block implicitly cites the Kerckhoffs threat model of
Section~\ref{sec:bg:lifecycle}: the attacker is assumed to know the
vault's design, parameters, and axis values. Where a field is
genuinely N/A for a given class, the cell reads ``N/A --- [reason]''
rather than being omitted.

\begin{table}[t]
\centering
\caption{Rationality and Scope matrix for the six attack classes.
Columns follow the seven-field schema of
Section~\ref{sec:bg:axes}: attacker capability (Cap.), defender
model (Def.), economic preconditions (Econ.), protocol/network
assumptions (Prot.), scope of validity (Scope), rational-attacker
check (Rat.), counter-factual mitigations (CF). Abbreviations: WT =
watchtower; RBF3 = BIP-125 Rule 3; \texttt{dust} = BIP-140/P2TR
$330$-sat or P2WPKH $546$-sat threshold.}
\label{tab:rationality}
\tiny
\setlength{\tabcolsep}{2pt}
\renewcommand{\arraystretch}{1.25}
\begin{tabular}{@{}p{0.9cm}p{1.55cm}p{1.55cm}p{1.55cm}p{1.55cm}p{1.55cm}p{1.55cm}p{1.55cm}@{}}
\toprule
Class & Cap. & Def. & Econ. & Prot. & Scope & Rat. & CF \\
\midrule
\classfee{} (TM1) &
  Hot key; anchor-desc. capital ($\approx 14{,}300$~sats). &
  Online WT; Core default mempool. &
  $r$-invariant; TM1 cost/$V$ $< 3\%$ at any $r$. &
  BIP-125 25-desc. cap; $3\satvB{}$ Core default dustRelay. &
  Applies iff \axfee{}\,=\,\texttt{out-of-band-anchor}; N/A for dynamic fee models. &
  Rational at any $r$ given attacker has hot key; yield $\approx V$. &
  TRUC/BIP-431 collapses desc.\,cap to 1 --- eliminates class; cluster-mempool: scope-limit.
\\
\classgrief{} (TM2) &
  \textbf{Zero keys} for \opccv{}; recovery-address public under Kerckhoffs. &
  Any; attack is pure liveness denial. &
  Linear in $r$; $\gamma(D)$ $r$-invariant. &
  Mempool-observable trigger tx; no RBF dependence. &
  Applies iff \axauth{}\,=\,\texttt{keyless}; reduces to TM6 otherwise. &
  Irrational without external incentive (no financial gain). &
  \axauth{}\,=\,\texttt{recoveryauth-key} (OPV) or key-gated (CTV/\catcsfs) eliminates class.
\\
\classscale{} (TM4) &
  Trigger key; capital for $N$ partial-unvaults. &
  WT with recovery authority; batching flag $b$; minimal-disruption policy maximises attacker leverage. &
  Rate-dep.; threshold $r^\dagger_{\mathrm{TM4}}$ per Eq.~\ref{eq:r-dagger}. &
  Dust thresholds; RBF3; sibling-eviction disallowed; attacker's revault chain is single-chain. &
  Applies iff \axrevault{}\,=\,\texttt{yes}; structural cost tail independent of adversary. &
  Rational only above $r^\dagger$ with external incentive; pure liveness attacker irrational at any $r$. &
  Defender batching: $r^\dagger$ $\times 1.8$--$2.7$; full-sweep-on-first-alert: 1 round only.
\\
\classhot{} (TM3) &
  Hot key; optional fee-pinning capital (if \classfee{} composes). &
  Online WT able to race during CSV; race bias toward attacker under TM1 composition. &
  Race outcome $r$-independent except via TM1. &
  CSV delay; mempool propagation (Ria23 cycling out of scope). &
  Applies iff \axwdbind{}\,$\neq$\,\texttt{signature-dual-bind}. &
  Rational only paired with fee-pinning budget; otherwise defender wins race generically. &
  \axwdbind{}\,=\,\texttt{signature-dual-bind} (CAT+CSFS): Schnorr EUF-CMA forces output-binding --- eliminates class.
\\
\classcold{} (TM11) &
  Cold key (Kerckhoffs). &
  Any; attack instant, no CSV window, no WT intervention. &
  $r$-invariant. &
  N/A --- pure signature check; no mempool dynamics. &
  Applies iff \axrecbind{}\,=\,\texttt{plain-signature}. &
  Rational whenever cold key is compromised. &
  \axrecbind{}\,=\,\texttt{signature-dual-bind} via Poelstra's \catcsfs{} recovery-leaf sketch: eliminates class.
\\
\classmode{} (TM8) &
  Any observer; no keys. &
  Developer discipline around mode-byte selection. &
  N/A --- full vault bypass at constant $110\vB{}$/spend. &
  BIP-341/342 \texttt{OP\_SUCCESSx} semantics: undefined mode terminates script with success. &
  Applies iff \axopsurf{}\,=\,\texttt{multi-mode-with-OP\_SUCCESS}. &
  Rational whenever a misconfigured vault exists on-chain. &
  \axopsurf{}\,=\,\texttt{single-mode}; bitmask-gated mode check; pymatt \texttt{CCV\_FLAG\_*} constants.
\\
\bottomrule
\end{tabular}
\end{table}

\providecommand{\todo}[2][]{\textbf{[TODO: #2]}}

% Appendix C: Per-experiment Measurements

\section{Per-Experiment Measurements}
\label{app:measurements}

This appendix collects the regtest measurements underlying
Section~\ref{sec:empirical} at greater granularity than the main
body's summary tables permit. Each value is structurally
deterministic (range~=~0 across repeated runs on the reference
implementations: \texttt{simple-ctv-vault}, \texttt{pymatt},
\texttt{opvault-demo}, our \texttt{cat-csfs-vault}, and our
\texttt{simple-simplicity-vault}).

\subsection{Revault Amplification}
\label{app:meas:revault}

Per-step vsize for nine successive partial withdrawals of
$5{,}000{,}000$~sats each from an initial $49{,}999{,}900$-sat
vault. Per-step size is constant across all nine steps for every
design, confirming structural determinism.

\begin{table}[h]
\centering
\caption{Per-step vsize for repeated partial unvaults with
revault. \opccv{} and \opvault{} measure their native
\texttt{trigger\_and\_revault}; \opctv{}, \catcsfs{}, and Simplicity do
not support native revault---the figures listed show the
single-withdrawal cost that repeats per step in a full-cycle
emulation (unvault $\rightarrow$ withdraw $\rightarrow$
re-deposit).}
\label{tab:revault_appendix}
\small
\begin{tabular}{@{}lrrl@{}}
\toprule
Design   & Per-step vB & Cumulative @ 9 steps (vB) & Revault model \\
\midrule
\opccv{}       & 111 & 999    & Native partial unvault \\
\opvault{} & 121 & 1{,}089 & Native partial unvault \\
\opctv{}       & 152 & 1{,}368 & Full-cycle emulation \\
\catcsfs{}  & 210 & 1{,}890 & Full-cycle emulation \\
Simplicity & 386 & 3{,}474 & Full-cycle emulation (Elements serialisation) \\
\bottomrule
\end{tabular}
\end{table}

\subsection{Batching Sweep}
\label{app:meas:batching}

Total transaction vsize as a function of batch size $N$ for a
single trigger aggregating $N$ vault UTXOs. \opccv{} and \opvault{}
amortise fixed overhead across inputs (sub-linear scaling); \opctv{},
\catcsfs{}, and Simplicity lack a spec-level batching primitive, so
the natural cost is $N$ independent transactions and the totals
grow strictly linearly.

\begin{table}[h]
\centering
\caption{Batched trigger vsize measurements at batch sizes
$N \in \{1, 2, 3, 5, 10, 20, 50\}$. Regression on the \opccv{} and
\opvault{} rows yields $\text{vsize}_{\text{trigger}}(N) =
\alpha_t + \beta_t \cdot N$ with
$(\alpha_t, \beta_t)_{\mathrm{OP\_CCV}} \approx (53, 100)$\vB{} and
$(\alpha_t, \beta_t)_{\opvault{}} \approx (198, 94)$\vB{}. The
recovery-side linear fits used in
Section~\ref{sec:empirical:batching} are
$(\alpha, \beta)_{\mathrm{OP\_CCV}} = (54, 68)$\vB{} and
$(\alpha, \beta)_{\opvault{}} = (154, 92)$\vB{}.}
\label{tab:batching_appendix}
\footnotesize
\setlength{\tabcolsep}{4pt}
\begin{tabular}{@{}lrrrrrrr@{}}
\toprule
Design    & $N{=}1$ & $N{=}2$ & $N{=}3$ & $N{=}5$ & $N{=}10$ & $N{=}20$ & $N{=}50$ \\
\midrule
\opccv{}       & 154 & 254  & 355   & 555    & 1{,}056  & 2{,}059  & 5{,}066 \\
\opvault{} & 292 & 385  & 479   & 667    & 1{,}135  & 2{,}073  & 4{,}885 \\
\midrule
\opctv{}       & 94  & 188  & 282   & 470    & 940      & 1{,}880  & 4{,}700 \\
\catcsfs{}  & 221 & 442  & 663   & 1{,}105 & 2{,}210 & 4{,}420  & 11{,}050 \\
Simplicity & 298 & 596 & 894   & 1{,}490 & 2{,}980 & 5{,}960  & 14{,}900 \\
\bottomrule
\end{tabular}
\end{table}

\subsection{Fee-Rate Sensitivity Projections}
\label{app:meas:fees}

Projected on-chain lifecycle cost
$C_{\text{lifecycle}}(r) = \text{vsize}_{\text{total}} \times r$
for each Bitcoin design at six historically observed fee rates.
Fee-pinning cost is $r$-invariant. Watchtower-exhaustion
feasibility is fee-dependent: the rightmost column gives the
number of per-UTXO splits an attacker needs to drive a $0.5$~BTC
\opccv{} vault below dust. The ranking of lifecycle totals is
fee-invariant because vsize ratios are constant.

\begin{table}[h]
\centering
\caption{Projected lifecycle cost (sats) across Bitcoin designs
and the \opccv{} split-to-dust count, at $V = 0.5$~BTC. Simplicity
runs on Elements and uses an incomparable fee mechanism (explicit
fee outputs), so it is excluded from Bitcoin-fee projections.}
\label{tab:fee_sensitivity_appendix}
\footnotesize
\begin{tabular}{@{}rrrrrrr@{}}
\toprule
& \multicolumn{4}{c}{Lifecycle cost (sats)} & Fee-pin & \opccv{} splits \\
\cmidrule(lr){2-5}
$r$ (\satvB{}) & \opctv{} & \opccv{} & \opvault{} & \catcsfs{} & (all $r$) & to dust \\
\midrule
1   &     368 &     565 &     567 &     553 & 14{,}300 & 409{,}836 \\
10  &   3{,}680 &   5{,}650 &   5{,}670 &   5{,}530 & 14{,}300 &  40{,}984 \\
50  &  18{,}400 &  28{,}250 &  28{,}350 &  27{,}650 & 14{,}300 &   8{,}197 \\
100 &  36{,}800 &  56{,}500 &  56{,}700 &  55{,}300 & 14{,}300 &   4{,}098 \\
300 & 110{,}400 & 169{,}500 & 170{,}100 & 165{,}900 & 14{,}300 &   1{,}366 \\
500 & 184{,}000 & 282{,}500 & 283{,}500 & 276{,}500 & 14{,}300 &     820 \\
\bottomrule
\end{tabular}
\end{table}

\subsection{Address Reuse}
\label{app:meas:reuse}

A second deposit to the same \opctv{} vault address creates a UTXO
whose amount does not match the committed template hash; the
second UTXO is permanently unspendable via the vault path
(confirmed on regtest). \opccv{}, \opvault{}, \catcsfs{}, and Simplicity
all handle repeated deposits safely: \opccv{} and \opvault{} validate
each UTXO per-instance against the contract; \catcsfs{} generates a
unique vault plan per deposit (distinct P2TR addresses);
Simplicity creates a fresh plan per instance. This is a
design-level property of \opctv{}'s input-count-and-amount commitment,
not an implementation artefact.

\subsection{Measurements Not Captured}
\label{app:meas:limits}

Regtest mines blocks on demand and has no fee market, so temporal
race dynamics (TM3 attacker-vs-watchtower timing, fee-pinning
propagation delays) cannot be measured on the harness. Structural
vsize is identical on regtest and mainnet for any fixed
transaction structure, so the lifecycle, batching, and revault
tables above transfer directly; the $r^\dagger$ thresholds of
Section~\ref{sec:empirical:batching} are derived by projection,
not timing.

\providecommand{\todo}[2][]{\textbf{[TODO: #2]}}

% Appendix D --- Simplicity on Elements: the dominant-corner existence proof.
% Content lifted from earlier §4.5 "Simplicity on Elements: A
% Cross-Substrate Reference" when §4 was compressed to main-body
% measurements only. Simplicity is not a Bitcoin covenant
% proposal and is excluded from the formal impossibility results in §3;
% it appears here as a constructive existence proof that the dominant
% corner D^\star is occupiable on some substrate.

\section{Simplicity on Elements: The Unoccupied Optimal Corner}
\label{app:simplicity}

\subsection{Scope and motivation}
\label{app:simplicity:scope}

Simplicity~\cite{OC17} is an \emph{alternative script primitive}---a
typed combinator calculus with Coq-verified denotational
semantics---deployed on Elements~\cite{DPNS16}, a Bitcoin-adjacent
sidechain whose blocks are produced by known functionaries rather
than by proof-of-work miners. Execution uses native-code
\emph{jets} for known combinator subtrees. Because Elements'
consensus model differs from Bitcoin's, we exclude Simplicity from
the formal results of Section~\ref{sec:impossibility} and report it
here as a cross-substrate reference: Simplicity demonstrates that
the dominant corner $D^\star$ of
Section~\ref{sec:imposs:design_space} is occupiable on some
substrate, even while no Bitcoin covenant proposal
occupies it.

This distinction is load-bearing. The claim of the paper is
\emph{unoccupied on Bitcoin}, not \emph{theoretically
unoccupiable}. Simplicity is the constructive existence proof for
the latter.

\subsection{Design-space position}
\label{app:simplicity:axes}

Our Simplicity vault occupies the strict-best axis values on every
axis where such a value exists:

\begin{itemize}
  \item \axauth{} = \texttt{cold-key}: key-gated recovery via a
    Schnorr signature check.
  \item \axrevault{} = \texttt{no}: atomic withdrawal; the program
    enforces the complete output tuple via
    \texttt{jet::outputs\_hash()\,=\,param::VAULT\_RECOVER\_OUTPUTS\_HASH}
    (see \texttt{simf/vault.simf}).
  \item \axfee{} = \texttt{dynamic} (Elements fee structure):
    fees live inside the transaction; no out-of-band anchor.
  \item \axwdbind{} = \texttt{dedicated-jets-combinator}: the
    withdrawal destination is bound through a typed output-hash
    commitment at compile time, stronger than any of the
    Bitcoin proposals on the same axis.
  \item \axrecbind{} = \texttt{dedicated-jets-combinator}: same
    binding for the recovery leaf, eliminating
    class \classcold{}.
  \item \axopsurf{} = \texttt{single-mode}: a program's meaning is
    determined by its Merkle root; no \texttt{OP\_SUCCESSx}-style
    fallthrough surface.
  \item \axintro{} = \texttt{typed-combinator-introspection}: the
    fundamental operators are total, typed, and Coq-verified.
\end{itemize}

Consequently, Simplicity's vault program is immune by construction
to classes \classfee{}, \classgrief{}, \classscale{}, \classhot{},
\classcold{}, and \classmode{} --- every attack class identified in
Section~\ref{sec:impossibility}. It occupies the strict-best
position on \axauth{}, \axrevault{}, \axwdbind{}, \axrecbind{}, and
\axintro{} simultaneously; no Bitcoin proposal does.

\subsection{Lifecycle measurements}
\label{app:simplicity:empirical}

The Simplicity vault lifecycle totals 865\vB{} (deposit 269,
trigger 298, withdraw 298, recover 286), $49$--$135\%$ larger than
the Bitcoin designs of Table~\ref{tab:vsizes}. The overhead bundles
Elements' serialisation (${\sim}30$--$40$ bytes per output for
confidential-transaction fields even in non-confidential mode) and
Simplicity's bit-oriented DAG witness format. Elements functionaries
do not operate a public mempool, so class \classfee{} (fee-channel
pinning) has no substrate analogue under honest-federation
assumptions.

Our Simplicity vault is atomic: recovery binds the complete output
tuple via the jets commitment cited above, admitting neither partial
withdrawal nor batched recovery and placing the program at
\axrevault{} = \texttt{no} in the design-space lattice. Class
\classscale{} (per-UTXO recovery-cost scaling) has no surface. Cold-key
compromise (TM11) forces recovery but cannot redirect funds, matching
\opctv{} and \opvault{}. A different Simplicity program could express
\axrevault{} = \texttt{yes} by replacing the whole-tuple commitment
with per-output checks; that design would inherit \classscale{}
under D2 (Section~\ref{sec:imposs:class_scale}) but is not what we
measured.

\subsection{Foundations and substrate}
\label{app:simplicity:foundations}

Simplicity's Coq-verified denotational semantics give stronger
guarantees than Alloy bounded model checking of the Bitcoin Script
extensions can produce --- a Simplicity program's meaning is
determined by its Merkle root, under a formal semantics that has
been mechanically verified. However, Simplicity on Liquid is live
today under \emph{federation trust}, whereas the four Bitcoin Script
extensions remain under \emph{Bitcoin consensus review}. The
substrate choice is a security-relevant decision alongside the
per-covenant analysis of Section~\ref{sec:impossibility}: a vault
on Simplicity inherits the federation's liveness and censorship
properties; a vault on Bitcoin inherits proof-of-work's. The
dominant corner being occupiable on Simplicity is therefore not a
turn-key prescription to migrate vaults to Liquid; it is evidence
that the unoccupied status among current Bitcoin proposals reflects a
consensus-substrate constraint rather than a design impossibility.


\end{document}
